\documentclass[11pt,a4paper]{article}

\usepackage[no-math]{fontspec}
\usepackage{polyglossia}
\usepackage{indentfirst}
\setdefaultlanguage{thai}
\setotherlanguage{english}

\setmainfont{TH Sarabun New}[
    Scale          = 1.46,
    BoldFont       = TH Sarabun New Bold,
    ItalicFont     = TH Sarabun New Italic,
    BoldItalicFont = TH Sarabun New Bold Italic,
]
\setsansfont{TH Sarabun New}[
    Scale          = 1.46,
    BoldFont       = TH Sarabun New Bold,
    ItalicFont     = TH Sarabun New Italic,
    BoldItalicFont = TH Sarabun New Bold Italic,
]
\newfontfamily\thaifont[
    Script         = Thai,
    Scale          = 1.46,
    BoldFont       = TH Sarabun New Bold,
    ItalicFont     = TH Sarabun New Italic,
    BoldItalicFont = TH Sarabun New Bold Italic,
]{TH Sarabun New}
\setmonofont{Consolas}[Scale=MatchLowercase, AutoFakeBold=1.4, AutoFakeSlant=0.2]
\newfontfamily\thaifonttt{Consolas}[Scale=MatchLowercase, AutoFakeBold=1.4, AutoFakeSlant=0.2]

\usepackage[top=2.0cm,bottom=2.1cm,left=2.2cm,right=2.0cm,footskip=1.0cm]{geometry}
\usepackage{graphicx,longtable,array,booktabs,xcolor,colortbl,float,enumitem,needspace,caption}
\usepackage{amsmath,amssymb,mathtools}
\usepackage{unicode-math}
\setmathfont{latinmodern-math.otf}
\PassOptionsToPackage{hyphens}{url}
\usepackage{fancyhdr,titlesec,titletoc,listings,url,hyperref}
\def\UrlBreaks{\do\/\do\-\do\_\do\.\do\:}

\XeTeXlinebreaklocale "th"
\XeTeXlinebreakskip = 0pt plus 0.42pt

% กำหนด \normalsize และ \small ให้มีระยะบรรทัดคงที่ตลอดทุกสภาพแวดล้อม (11pt x 1.46 ≈ 16pt TH Sarabun New)
\makeatletter
\renewcommand\normalsize{%
   \@setfontsize\normalsize{11}{16.6}%
   \abovedisplayskip 6\p@
   \abovedisplayshortskip 6\p@
   \belowdisplayshortskip 6\p@
   \belowdisplayskip 6\p@
   \let\@listi\@listI}
\renewcommand\small{%
   \@setfontsize\small{10}{14.8}%
   \abovedisplayskip 5\p@
   \abovedisplayshortskip 5\p@
   \belowdisplayshortskip 5\p@
   \belowdisplayskip 5\p@}
\makeatother
\normalsize

% โทนสีเอกสารวิชาการ CPE: ตารางสีฟ้าอ่อน ศัพท์โค้ดสีเขียว และคำบรรยายภาพ/ตารางตัวเอียงสีฟ้า 11pt (อ้างอิงขนาดฟอนต์มาตรฐานรายงาน CPE232)
\definecolor{TableHead}{HTML}{E6EEF6}
\definecolor{TableStripe}{HTML}{F8FAFC}
\definecolor{RuleGray}{HTML}{C6CDD5}
\definecolor{DarkBorder}{HTML}{64748B}
\definecolor{CodeTint}{HTML}{F8FAFC}
\definecolor{KeyGreen}{HTML}{087A46}
\definecolor{CaptionBlue}{HTML}{0000EE}
\definecolor{LinkNavy}{HTML}{1E3A8A}

\hypersetup{
    unicode=true,
    colorlinks=true,
    linkcolor=black,
    urlcolor=KeyGreen,
    citecolor=black,
    pdftitle={CPE342 Term Project: Karena Player Intelligence System},
    pdfauthor={Group ML004},
    hypertexnames=false
}
\urlstyle{same}
\setlength{\emergencystretch}{1.1em}
\widowpenalty=10000
\clubpenalty=10000
\setlength{\parindent}{1.0cm}
\setlength{\parskip}{4.5pt}
\setlength{\tabcolsep}{5.5pt}
\renewcommand{\arraystretch}{1.26}
\newcolumntype{L}[1]{>{\raggedright\arraybackslash}p{#1}}
\setlist[itemize]{leftmargin=1.9em,itemsep=2pt,topsep=0pt,partopsep=0pt,parsep=2pt}
\setlist[enumerate]{leftmargin=1.9em,itemsep=2pt,topsep=0pt,partopsep=0pt,parsep=2pt}

\setcounter{tocdepth}{2}
\setcounter{secnumdepth}{3}

% ขนาดหัวข้อและเนื้อหาอ้างอิงมาตรฐานรายงาน (TH Sarabun New Scale 1.46: บทที่ 14pt, หัวข้อรอง 13pt, ย่อย 11.5pt, เนื้อหา 11pt≈16pt)
\titleformat{\section}{\normalfont\raggedright\color{black}\fontsize{14}{21}\selectfont\bfseries}{บทที่~\thesection\quad}{0pt}{}[\vspace{4pt}{\color{RuleGray}\hrule height 0.7pt}]
\titleformat{name=\section,numberless}{\normalfont\raggedright\color{black}\fontsize{14}{21}\selectfont\bfseries}{}{0pt}{}[\vspace{4pt}{\color{RuleGray}\hrule height 0.7pt}]
\titleformat{\subsection}{\normalfont\raggedright\color{black}\fontsize{13}{19.5}\selectfont\bfseries}{\thesubsection\quad}{0pt}{}
\titleformat{\subsubsection}{\normalfont\raggedright\color{black}\fontsize{11.5}{17.5}\selectfont\bfseries}{\thesubsubsection\quad}{0pt}{}
\titlespacing{\section}{0pt}{3pt}{7pt}
\titlespacing{\subsection}{0pt}{9pt}{3.5pt}
\titlespacing{\subsubsection}{0pt}{6pt}{2.5pt}

% สารบัญแบบเป็นทางการ (ไร้จุดไข่ปลา ขนาด 11pt เท่าเนื้อหา และจัดเลขหน้าชิดขอบขวา \textwidth พอดี)
\contentsmargin{1.4em}
\titlecontents{section}[0em]{\filright\addvspace{4.5pt}\fontsize{11}{16.2}\selectfont\bfseries}{\mbox{บทที่~\thecontentslabel}\quad}{}{\hfill\contentspage}
\titlecontents{subsection}[2.2em]{\filright\addvspace{1pt}\fontsize{11}{16.2}\selectfont}{\thecontentslabel\quad}{}{\hfill\contentspage}
\titlecontents{subsubsection}[4.2em]{\filright\addvspace{0.5pt}\fontsize{10.5}{15.4}\selectfont}{\thecontentslabel\quad}{}{\hfill\contentspage}

% คำบรรยายภาพและตารางขนาด 11pt ตัวเอียงสีฟ้า จัดกึ่งกลางมาตรฐาน
\DeclareCaptionFont{captext}{\color{CaptionBlue}\itshape\fontsize{11}{16.2}\selectfont}
\captionsetup{font=captext,labelfont=bf,justification=centering,singlelinecheck=true,labelsep=space,skip=6pt}

% คำสั่งแสดงชื่อฟีเจอร์/ไฟล์/พารามิเตอร์ด้วยอักษรโค้ดสีเขียวและขนาดเท่าข้อความรอบข้าง
\DeclareRobustCommand{\code}[1]{\leavevmode{\ttfamily\color{KeyGreen!85!black}\let\codeunderscore\_\def\_{\codeunderscore\allowbreak}\hyphenpenalty=50\exhyphenpenalty=50\relax#1}}
\newcommand{\wideimg}[2][1.06\textwidth]{\makebox[\textwidth][c]{\includegraphics[width=#1]{#2}}}
\newcommand{\tblfont}{\fontsize{10}{15.0}\selectfont\arrayrulecolor{RuleGray}}

\lstset{
    basicstyle=\fontspec{Consolas}\fontsize{9.5}{12.5}\selectfont,
    columns=fullflexible,
    keepspaces=true,
    breaklines=true,
    breakatwhitespace=false,
    showstringspaces=false,
    backgroundcolor=\color{CodeTint},
    frame=single,
    rulecolor=\color{DarkBorder},
    framesep=6pt,
    aboveskip=6pt,
    belowskip=7pt,
    xleftmargin=.15cm,
    tabsize=4,
    keywordstyle=\color{blue!75!black}\bfseries,
    commentstyle=\color{gray}\itshape,
    stringstyle=\color{teal!70!black}
}

% ไม่มีหัวกระดาษ (Zero Running Header) แสดงเฉพาะเลขหน้ามุมขวาล่าง
\pagestyle{fancy}
\fancyhf{}
\fancyfoot[R]{\fontsize{11}{13}\selectfont\thepage}
\renewcommand{\headrulewidth}{0pt}
\renewcommand{\footrulewidth}{0pt}
\fancypagestyle{plain}{
    \fancyhf{}
    \fancyfoot[R]{\fontsize{11}{13}\selectfont\thepage}
    \renewcommand{\headrulewidth}{0pt}
    \renewcommand{\footrulewidth}{0pt}
}

\newcommand{\thaiFrontPage}{%
    \ifcase\value{page}\or ก\or ข\or ค\or ง\or จ\or ฉ\or ช\or ซ\else\arabic{page}\fi%
}

\AtBeginDocument{%
    \normalsize
    \lineskiplimit=-3pt
    \lineskip=1pt
    \tolerance=1200\emergencystretch=1.2em\hbadness=2500
    \renewcommand{\contentsname}{สารบัญ}
    \renewcommand{\tablename}{ตารางที่}
    \renewcommand{\figurename}{รูปที่}
    \renewcommand{\refname}{บรรณานุกรม}
}

\begin{document}

% หน้าปก
\newgeometry{top=1.8cm,bottom=2.2cm,left=2.5cm,right=2.5cm}
\InputIfFileExists{frontpage/frontpage.tex}{}{\typeout{frontpage/frontpage.tex not found}}
\restoregeometry

% สารบัญ (เริ่มต้นทันทีหลังหน้าปก ด้วยเลขหน้าอักษรไทย ก, ข...)
\clearpage
\pagenumbering{arabic}
\renewcommand{\thepage}{\thaiFrontPage}
\renewcommand{\contentsname}{สารบัญ}
\tableofcontents

% เนื้อหาหลัก: เลขอารบิกเริ่มที่ 1 (จัดกระจายข้อความชิดขอบซ้าย-ขวาอย่างสมดุล พร้อมย่อหน้าบรรทัดแรก 1.0cm)
\clearpage
\pagenumbering{arabic}
\renewcommand{\thepage}{\arabic{page}}
\setlength{\parindent}{1.0cm}
\hyphenpenalty=50
\exhyphenpenalty=50

% ==============================================================================
\section{บทนำ}

\subsection{ที่มาและความสำคัญ}

ในปัจจุบัน อุตสาหกรรมเกมออนไลน์และแพลตฟอร์มความบันเทิงดิจิทัลมีการแข่งขันสูงขึ้นอย่างต่อเนื่อง ผู้ให้บริการมิได้แข่งขันกันเฉพาะคุณภาพของตัวเกมเท่านั้น แต่ยังต้องอาศัยการวิเคราะห์ข้อมูลขนาดใหญ่และเทคโนโลยี Machine Learning เพื่อรักษาความยุติธรรมในการแข่งขัน เข้าใจพฤติกรรมผู้เล่น เพิ่มประสิทธิภาพการสร้างรายได้ และบริหารจัดการมูลค่าตลอดอายุการใช้งานของผู้เล่น (Player Lifetime Value: LTV) อย่างเป็นระบบ

การแข่งขัน \textbf{Karena Player Intelligence System} จำลองบทบาททีม \textbf{Machine Learning Engineer} ของบริษัท \textbf{Karena Thailand} ผู้พัฒนาและให้บริการแพลตฟอร์มเกมออนไลน์ขนาดใหญ่ในภูมิภาค (ลักษณะเดียวกับแพลตฟอร์ม \textit{Garena}) ซึ่งดูแลเกมหลัก 5 เกมที่มีแนวการเล่นแตกต่างกัน ได้แก่ \textit{PUBE Mobile} (เกมยิงเอาชีวิตรอดเชิงกลยุทธ์), \textit{Free Fried} (เกมยิงแข่งขันจังหวะเร็ว), \textit{7-11 Knights} (เกม RPG สะสมตัวละคร), \textit{Roblock} (เกมสร้างโลกจำลองผู้เล่นหลายคน) และ \textit{FiveN} (เกมชุมชนจำลองบทบาทสมมติ) จากการเติบโตของฐานผู้เล่นอย่างรวดเร็ว บริษัทเผชิญปัญหาเชิงปฏิบัติการสำคัญ 5 ด้านที่ส่งผลกระทบโดยตรงต่อประสบการณ์ผู้เล่นและผลประกอบการ ดังนี้

\begin{enumerate}[leftmargin=1.9em,itemsep=2.5pt]
    \item \textbf{ปัญหาภาระงานคัดกรองการโกง (Task 1 — Anti-Cheat Pre-Filter):} แพลตฟอร์มได้รับรายงานการโกงมากกว่า 10,000 กรณีต่อวัน ขณะที่เจ้าหน้าที่ผู้ตรวจสอบสามารถพิจารณาด้วยตนเองได้เพียง 500 กรณีต่อวัน ก่อให้เกิดงานค้างสะสมกว่า 45,000 กรณี ระบบกฎเกณฑ์อัตโนมัติแบบเดิมมีอัตรา False Positive สูงถึง 70\% ทำให้ผู้เล่นสุจริตถูกระบุสถานะผิดพลาดจนนำไปสู่การระงับบัญชีโดยไม่สมควรและสร้างภาระแก่ฝ่ายบริการลูกค้า ขณะเดียวกันผู้เล่นทุจริตที่ใช้โปรแกรม Subtle Aimbot (ช่วยเล็งอัตโนมัติ) และ Wallhack (มองทะลุสิ่งกีดขวาง) ยังคงหลบเลี่ยงการตรวจจับได้ ส่งผลให้อัตราการรักษาผู้เล่นในโหมดแข่งขันลดลง 15\% และกระทบต่อรายได้ราว 15 ล้านบาทต่อวัน
    \item \textbf{ปัญหาการตลาดแบบไม่ตรงกลุ่มเป้าหมาย (Task 2 — Player Segment Classification):} ฝ่ายการตลาดใช้งบประมาณกว่า 60 ล้านบาทต่อเดือนไปกับแคมเปญการตลาดแบบไม่จำแนกกลุ่ม (Mass Campaign) ซึ่งมีอัตราการตอบรับเพียง 2--5\% เช่น การแสดงประกาศแข่งขัน Esports ต่อผู้เล่นสาย Casual บนมือถือ การส่งโฆษณาจำหน่ายไอเทมเครื่องแต่งกาย (Cosmetic) ต่อผู้ที่ไม่เคยเติมเงิน หรือการส่งคำเชิญกิจกรรมสวมบทบาทของเกม \textit{FiveN} ให้กลุ่มผู้เล่นสายแข่งขันของ \textit{PUBE Mobile} ทั้งที่ผู้เล่นกลุ่มใช้จ่ายสูง (Whales) เพียง 5\% แรกสร้างรายได้ถึง 60\% ของแพลตฟอร์ม และตัวชี้วัดอย่างชั่วโมงการเล่นเพียงอย่างเดียวไม่อาจแยกผู้เล่นสายไต่แรงก์ (Grinder), สายสังคม (Social) และสายเติมหนัก (Whale) ออกจากกันได้ การจำแนกกลุ่มอย่างแม่นยำจึงมีศักยภาพเพิ่มประสิทธิผลแคมเปญได้ถึง 3--4 เท่า
    \item \textbf{ความคลาดเคลื่อนในการพยากรณ์รายได้และดูแลลูกค้า VIP (Task 3 — Monthly Spending Prediction):} ฝ่ายการเงินคาดการณ์รายได้รายไตรมาสคลาดเคลื่อน 35--40\% จนกระทบการวางงบประมาณและการจัดสรรบุคลากร เนื่องจากยอดเติมเงินมีลักษณะ Zero-Inflated และมีหางขวาหนัก (Heavy Right Tail) ขณะเดียวกันทีม VIP Support จำนวน 20 คนขาดข้อมูลคาดการณ์รายบุคคล จึงทุ่มทรัพยากรดูแลผู้เล่นที่เข้าใช้งานต่อเนื่องแต่มียอดใช้จ่ายต่ำ 150 รายซึ่งมียอดใช้จ่ายรวมเพียง 12,000 บาท แต่กลับปล่อยให้ลูกค้า Whale ที่ใช้จ่ายถึง 255,000 บาทต้องรอตอบคำถามนานถึง 3 วันจนย้ายไปแพลตฟอร์มคู่แข่ง นอกจากนี้การแจกโค้ดส่วนลด 50\% แบบสุ่มให้กลุ่มที่พร้อมจ่ายเต็มราคาอยู่แล้วยังสร้างความสูญเสียรายได้กว่า 6 ล้านบาทต่อเดือน และทำให้กำหนดช่วงเวลากิจกรรมจำกัดเวลา (Limited-Time Events) ไม่สอดคล้องกับช่วงเวลาจับจ่ายของผู้เล่นหลัก
    \item \textbf{คอขวดในการคัดแยกภาพหน้าจอแจ้งปัญหา (Task 4 — Game Title Detection):} ฝ่ายบริการลูกค้าได้รับ Ticket แจ้งปัญหาพร้อมภาพหน้าจอที่ไม่มีป้ายกำกับราว 2,000 ใบต่อวัน โดยเจ้าหน้าที่ต้องใช้เวลามากกว่า 10 นาทีต่อใบในการระบุชื่อเกมก่อนส่งต่อผู้เชี่ยวชาญ อีกทั้งกิจกรรมชิงรางวัล ``Play Multiple Games for Rewards'' มีภาพส่งเข้าร่วมกว่า 50,000 ภาพ แต่ตรวจสอบด้วยมือได้เพียง 3\% จึงเกิดช่องโหว่การทุจริตรับรางวัล รวมถึงภาพจากการถ่ายทอดสด ฟอรัม และชุมชนผู้เล่นกว่า 100,000 ภาพต่อสัปดาห์ที่ฝ่ายการตลาดยังไม่สามารถจัดหมวดหมู่เพื่อติดตามความนิยมรายเกมแบบอัตโนมัติได้
    \item \textbf{การโจรกรรมบัญชีและการเอาเปรียบเศรษฐกิจในเกม (Task 5 — Account Security Monitoring):} แพลตฟอร์มพบเหตุบัญชีที่ถูกโจรกรรมและยืนยันแล้วกว่า 5,000 บัญชีต่อเดือน (ซึ่งคาดว่าปริมาณจริงสูงเป็น 2 เท่า) จนกลายเป็นปัญหาอันดับหนึ่งที่ใช้ทรัพยากรฝ่ายบริการลูกค้าถึง 60\% ควบคู่กับขบวนการรับจ้างปรับเพิ่มระดับแรงก์ (Account Boosting) การใช้งานบัญชีอัตโนมัติ (Bot Farming) ตลอด 24 ชั่วโมงที่ก่อให้เกิดภาวะเงินเฟ้อในเกม และการใช้ VPN เพื่อซื้อสินค้าในราคาภูมิภาคอื่น (ส่วนลด 70\%) ระบบกฎเกณฑ์แบบคงที่ (Static Rule-Based System) สร้างสัญญาณเตือนผิดพลาด (False Alarms) สูงถึง 60\% ต่อผู้เล่นที่เดินทางข้ามพื้นที่จริงหรือผู้เล่นที่พัฒนาฝีมือขึ้นตามธรรมชาติ ขณะที่ผู้โจมตีซึ่งเว้นช่วงเวลาทำรายการเกิน 24 ชั่วโมงสามารถหลบเลี่ยงการตรวจจับได้โดยง่าย
\end{enumerate}

\subsection{วัตถุประสงค์ของโครงงาน}

โครงงานนี้มีวัตถุประสงค์หลักเพื่อพัฒนาระบบต้นแบบสำหรับการแข่งขัน \textbf{Karena Player Intelligence System} ที่เชื่อมโยงแบบจำลอง Machine Learning ทั้ง 5 งานเข้าสู่ระบบคาดการณ์ \textbf{Player Lifetime Value (LTV)} ของแพลตฟอร์ม เพื่อตอบโจทย์เชิงยุทธศาสตร์ 4 ด้าน ได้แก่ การรักษาความยุติธรรมและความปลอดภัยของระบบ (Maintain Game Integrity), การทำความเข้าใจพฤติกรรมและความต้องการของผู้เล่น (Understand Player Behavior), การเพิ่มประสิทธิภาพกลยุทธ์การสร้างรายได้ (Optimize Monetization Strategies) และการยกระดับอัตราการรักษาผู้เล่นควบคู่กับมูลค่าตลอดอายุการใช้งาน (Improve Retention and LTV) โดยมีเป้าหมายเชิงเทคนิคดังนี้

\begin{enumerate}[leftmargin=1.8em,itemsep=2pt]
    \item เพื่อวิเคราะห์โครงสร้างข้อมูลเชิงสำรวจ (Exploratory Data Analysis: EDA) ตัวแปรนำเข้าตาม Data Card และกลไกการกำเนิดข้อมูล (Data-Generating Process) ของชุดข้อมูลทั้ง 5 งาน ทั้งในรูปแบบตารางพฤติกรรม อนุกรมเวลาหลายช่วง และภาพหน้าจอการเล่นเกม
    \item เพื่อออกแบบกระบวนการ Data Preprocessing, Feature Engineering เฉพาะโดเมน และสถาปัตยกรรมโมเดลที่ Train จากข้อมูลของการแข่งขันอย่างถูกต้องตามกติกา ปราศจาก Data Leakage และไม่ใช้เครื่องมือ AutoML
    \item เพื่อทำ Probability Calibration และจูน Decision Threshold ให้สอดคล้องกับตัวชี้วัดเฉพาะของแต่ละงาน ($F_2$, Macro-$F_1$, NMAE และ $F_3$) พร้อมพิสูจน์ผลด้วย Out-of-Fold Cross-Validation และ Paired Bootstrap Confidence Interval
    \item เพื่อบูรณาการผลลัพธ์จากทั้ง 5 งานเข้าสู่ไฟล์คำตอบรวม \code{submission.csv} ที่มีเสถียรภาพ วิเคราะห์ข้อค้นพบ ข้อผิดพลาดที่พบบ่อย และเชื่อมโยงสู่ผลกระทบเชิงธุรกิจของบริษัท Karena Thailand
\end{enumerate}

\subsection{ขอบเขต ข้อกำหนดการแข่งขัน และโครงสร้างผลลัพธ์ที่ต้องส่ง}

การพัฒนาโมเดลทั้งหมดดำเนินการภายใต้กรอบเวลาและกติกาของรายวิชา CPE 342 ภาคเรียนที่ 1 ปีการศึกษา 2569 โดยแบ่งช่วงการแข่งขันบน Kaggle ระหว่างวันที่ 2--24 พฤศจิกายน 2569 และกำหนดส่งรายงานทางวิชาการพร้อมคลังซอร์สโค้ดบน GitHub ผ่านระบบ LEB2 ภายในวันที่ 30 พฤศจิกายน 2569 สัดส่วนการประเมินประกอบด้วย \textbf{คะแนนบน Kaggle Leaderboard 70\%}, \textbf{รายงานทางวิชาการ (Technical Report) 20\%} และ \textbf{ซอร์สโค้ดและการทำซ้ำ (Source Code) 10\%}

ข้อกำหนดสำคัญของการแข่งขันที่ทีมผู้จัดทำปฏิบัติตามอย่างเคร่งครัดมี 4 ประการ ดังนี้
\begin{itemize}
    \item \textbf{การใช้ชุดข้อมูลภายในโครงการ:} ใช้เฉพาะชุดข้อมูลที่กำหนดให้เท่านั้น โดยไม่ดึงข้อมูลภายนอกหรือเก็บข้อมูลเพิ่มเติม
    \item \textbf{การป้องกันการใช้ข้อมูลทดสอบ (Test Set Isolation):} ห้ามนำข้อมูลชุดทดสอบ (\code{test.csv}) เข้าสู่การ Train, การจูน Hyperparameter หรือการทำ Supervised Pseudo-Labeling ใน Tasks 1--4 ยกเว้น Task 5 ที่ออกแบบเป็น Unsupervised Anomaly Detection บนชุดทดสอบโดยเฉพาะ
    \item \textbf{ขอบเขตการ Train โมเดล:} Tasks 1, 2, 3 และ 5 ต้อง Train โมเดลจากศูนย์ (Train from Scratch) ส่วน Task 4 อนุญาตให้ใช้ Pre-trained Vision Backbones ได้
    \item \textbf{ข้อจำกัดเครื่องมือ AutoML:} ห้ามใช้ AutoML Frameworks (เช่น AutoGluon, PyCaret, H2O AutoML, FLAML, TPOT, Auto-sklearn) แต่อนุญาตให้ใช้ไลบรารีจูน Hyperparameter เช่น Optuna และ Scikit-learn Search รวมถึงการสร้าง Ensemble ภายใต้กติกาข้างต้น
\end{itemize}

นอกจากนี้ จากการตรวจสอบไฟล์ข้อมูลจริงก่อนเริ่มพัฒนา พบประเด็นทางเทคนิคสำคัญ 2 ประการที่กำหนดเป็นมาตรฐานร่วมกันตลอด Pipeline ได้แก่
\begin{itemize}
    \item \textbf{มาตรฐาน Target Label ของ Task 4:} เอกสารสไลด์โจทย์หน้า 24 แสดงตัวอย่างคอลัมน์ \code{task4} เป็นข้อความชื่อเกม แต่ไฟล์ \code{task4/data/datacard.txt} และชุดข้อมูลจริงกำหนดใช้รหัสตัวเลขจำนวนเต็ม \code{0..4} (\code{0 = 7-11k}, \code{1 = fiveN}, \code{2 = freefried}, \code{3 = pube}, \code{4 = roblock}) คณะผู้จัดทำจึงยึดรหัส \code{0..4} เป็นมาตรฐานเดียวทั้งระบบ
    \item \textbf{ค่าน้ำหนักคะแนนรวม:} ค่าน้ำหนัก $w_1 \dots w_5$ สำหรับคำนวณคะแนนรวมบน Kaggle ไม่ได้เปิดเผยเป็นตัวเลข จึงมุ่งยกระดับประสิทธิภาพของแต่ละ Task ตามตัวชี้วัดทางการของงานนั้นโดยตรง
\end{itemize}

ช่องทางอ้างอิงการแข่งขันและคลังจัดเก็บไฟล์กลางของกลุ่ม \textbf{ML004} (Source Code \& Report Archive) ประกอบด้วย
{\raggedright
\begin{itemize}[leftmargin=1.8em,itemsep=1pt,topsep=0pt]
    \item \textbf{Kaggle:} \url{https://www.kaggle.com/competitions/player-intelligence-system-cpe-342-ml/}
    \item \textbf{คลังซอร์สโค้ดกลาง:} \url{https://github.com/beambeambeam/ml00x} (\code{task1}--\code{task5}, \code{merge.py})
    \item \textbf{คลังรายงานวิชาการ:} \url{https://github.com/beambeambeam/ml00x/tree/main/report} (\code{.tex/.pdf})
\end{itemize}\par}

\subsection{เครื่องมือและเทคโนโลยีที่ใช้ในการพัฒนา}

ทุกขั้นตอนของการวิเคราะห์ การทดลอง และการสร้างผลลัพธ์พัฒนาด้วยภาษา Python บน Jupyter Notebook ที่สามารถประมวลผลซ้ำ (Reproduce) ได้ครบถ้วน โดยมีเครื่องมือและไลบรารีหลักดังนี้
\begin{itemize}[itemsep=1.5pt,topsep=2pt]
    \item \textbf{การจัดการข้อมูลและการคำนวณเชิงตัวเลข:} NumPy, Pandas และ SciPy (สำหรับ Optimization ด้วยวิธี SLSQP, Gauss-Legendre Quadrature และ Statistical Testing)
    \item \textbf{Tabular Machine Learning \& Deep Tabular Models:} Scikit-learn, LightGBM, XGBoost, CatBoost และ PyTorch สำหรับ Neural Network บนข้อมูลตาราง
    \item \textbf{Computer Vision \& Pre-trained Backbones (Task 4):} PyTorch ร่วมกับคลังโมเดล \code{timm} และ Transformers สำหรับสถาปัตยกรรม ConvNeXt, ViT (EVA-02, DINOv2, DINOv3) และ VLM (SigLIP2, CLIP) บน GPU T4
    \item \textbf{การปรับแต่ง การแสดงผล และการควบคุมเวอร์ชัน:} Optuna สำหรับ Hyperparameter Tuning ภายใต้ Nested Cross-Validation, Matplotlib / Seaborn สำหรับสร้างกราฟวิเคราะห์ความละเอียด 300~dpi และ Git / GitHub สำหรับบริหารจัดการซอร์สโค้ดแบบแยกโมดูลระหว่างสมาชิกกลุ่ม \textbf{ML004}
\end{itemize}

% ==============================================================================
\clearpage
\section{กรอบการประเมินผล เกณฑ์พื้นฐาน และแนวทางป้องกัน Data Leakage}

\subsection{นิยามตัวชี้วัดทั้ง 4 ตระกูลและเหตุผลเชิงธุรกิจ}

แบบจำลองทั้ง 5 งานถูกประเมินด้วยตัวชี้วัดที่แตกต่างกันตามต้นทุนความผิดพลาดจริงทางธุรกิจ การเข้าใจโครงสร้างทางคณิตศาสตร์ของตัวชี้วัดจึงเป็นปัจจัยสำคัญที่กำหนดทั้ง Loss Function และการเลือก Decision Threshold

\begin{enumerate}
    \item \textbf{ตัวชี้วัด $F_\beta$ Score แบบไม่สมมาตร (Task 1 และ Task 5):}
    \[
    F_\beta = (1 + \beta^2)\frac{\mathrm{Precision} \cdot \mathrm{Recall}}{\beta^2 \cdot \mathrm{Precision} + \mathrm{Recall}} = \frac{(1 + \beta^2)\,\mathrm{TP}}{(1 + \beta^2)\,\mathrm{TP} + \beta^2\,\mathrm{FN} + \mathrm{FP}}
    \]
    \begin{itemize}
        \item \textbf{Task 1 ($F_2$ Score, $\beta=2$):} ให้น้ำหนัก False Negative ($\mathrm{FN}$) มากกว่า False Positive ($\mathrm{FP}$) ถึง \textbf{4 เท่า} (Recall สำคัญเป็น 2 เท่าของ Precision) เนื่องจากผู้เล่นทุจริตที่เล็ดลอดการตรวจจับก่อให้เกิดความเสียหายต่อผู้เล่นจำนวนมากในระบบ การคัดกรองเบื้องต้นเพื่อส่งต่อให้เจ้าหน้าที่ตรวจสอบจึงยอมรับผู้ต้องสงสัยเพิ่มได้มากกว่าการปล่อยให้ผู้เล่นทุจริตหลุดพ้นจากการตรวจสอบ
        \item \textbf{Task 5 ($F_3$ Score, $\beta=3$):} ให้น้ำหนัก $\mathrm{FN}$ มากกว่า $\mathrm{FP}$ ถึง \textbf{9 เท่า} (Recall สำคัญเป็น 3 เท่าของ Precision) เพื่อไม่ให้บัญชีที่ถูกยึดครองหรือบัญชีอัตโนมัติหลุดพ้นจากการตรวจจับของระบบเฝ้าระวังความปลอดภัย
    \end{itemize}

    \item \textbf{Unweighted Macro $F_1$ (Task 2 และ Task 4):}
    \[
    \mathrm{Macro\ } F_1 = \frac{1}{C}\sum_{c=0}^{C-1} \frac{2 \cdot \mathrm{Precision}_c \cdot \mathrm{Recall}_c}{\mathrm{Precision}_c + \mathrm{Recall}_c}
    \]
    ให้ความสำคัญทุกคลาสเท่ากันโดยไม่ขึ้นกับสัดส่วนคลาส (Class Proportion) ($C=4$ ใน Task 2 และ $C=5$ ใน Task 4) ดังนั้นคลาสขนาดเล็ก เช่น ผู้เล่นกลุ่ม Whale หรือเกมที่มีภาพตัวอย่างน้อย มีน้ำหนักต่อคะแนนเฉลี่ยเท่ากับคลาสขนาดใหญ่ การทำนายด้วยค่า \code{argmax} ปกติจากความน่าจะเป็นตั้งต้น (Raw Probabilities) จึงไม่เหมาะสมเมื่อคลาสไม่สมดุล

    \item \textbf{Normalized Mean Absolute Error Score (Task 3):}
    \[
    \mathrm{NMAE\ Score} = \frac{1}{1 + \dfrac{\sum_{i=1}^{N} |y_i - \hat{y}_i|}{\sum_{i=1}^{N} y_i}} \in [0, 1]
    \]
    โดยค่า 1.0 หมายถึงปราศจากความคลาดเคลื่อนสัมบูรณ์ เนื่องจาก $\smash{\sum_{i=1}^N} y_i$ เป็นค่าคงที่ของชุดข้อมูล การเพิ่มค่า NMAE Score ให้สูงสุดจึงสมมูลเชิงคณิตศาสตร์กับการลดผลรวมความคลาดเคลื่อนสัมบูรณ์ ($L_1$ Loss) ซึ่งคำตอบที่เหมาะสมที่สุดทางทฤษฎีคือ \textbf{ค่ามัธยฐานแบบมีเงื่อนไข (Conditional Median)} มิใช่ค่าเฉลี่ยคาดหมายแบบ $L_2$

    \item \textbf{คะแนนประเมินรวมระบบ (Weighted Score):}
    \[
    \mathrm{Weighted\ Score} = \frac{\sum_{i=1}^{5} w_i S_i}{\sum_{i=1}^{5} w_i} \in [0, 1]
    \]
    ใช้รวมผลคะแนน $S_1 \dots S_5$ จากทั้ง 5 งานเข้าเป็นดัชนีประเมินประสิทธิภาพรวมบน Kaggle Leaderboard
\end{enumerate}

\subsection{เกณฑ์คะแนนพื้นฐาน โครงสร้างชุดข้อมูล และรูปแบบการส่งคำตอบ}

สำหรับการประเมินผลในส่วนการแข่งขันบน Kaggle (70\%) ระบบกำหนดเกณฑ์อ้างอิงอย่างเป็นทางการเพียงเกณฑ์เดียว คือ \textbf{Kaggle Baseline (0.72601)} ซึ่งเป็นชุดคำตอบมาตรฐานจากผู้จัดรายวิชาบน Leaderboard เพื่อใช้วัดผลการพัฒนาและเป็นจุดแบ่งเกณฑ์คำนวณคะแนนเกรดผ่านการแจกแจง T-Score ของกลุ่มที่ทำคะแนนผ่านเกณฑ์ ($\mathrm{Score} \ge \mathrm{Baseline}$) ตามสมการ
\begingroup
\setlength{\abovedisplayskip}{3pt}\setlength{\belowdisplayskip}{3pt}%
\[
\mathrm{Grade\%} = \alpha + (100\% - \alpha)\,\Phi(t), \qquad t = \frac{\mathrm{Your\ Score} - \mu}{\sigma}
\]
\endgroup
\noindent โดยที่ $\mu$ และ $\sigma$ คือค่าเฉลี่ยและส่วนเบี่ยงเบนมาตรฐานของคะแนนกลุ่มที่สูงกว่าเกณฑ์ \textbf{Kaggle Baseline (0.72601)} ทั้งหมด ($\Phi(t)$ คือฟังก์ชันการแจกแจงสะสมแบบปกติมาตรฐาน) หากคะแนนต่ำกว่าเกณฑ์จะคิดคะแนนตามสัดส่วน $\alpha \cdot (\mathrm{Your\ Score} / \mathrm{Baseline})$ สำหรับไฟล์ส่งคำตอบรวม \code{sample\_submission.csv} กำหนดขนาด $25{,}889 \times 6$ คอลัมน์ (รหัส \code{id} ตั้งแต่ \code{ANS00001} ถึง \code{ANS25889} และ \code{task1}--\code{task5}) ดังสรุปภาพรวมทั้ง 5 งานใน \textbf{ตารางที่~\ref{tab:dataset_overview}}

\begin{table}[H]\centering
{\tblfont\renewcommand{\arraystretch}{0.95}\setlength{\tabcolsep}{4pt}
\caption{สรุปลักษณะงาน ขนาดชุดข้อมูล ตัวแปรเป้าหมาย ตัวชี้วัด และวิธีแบ่ง Validation ทั้ง 5 งาน}
\label{tab:dataset_overview}
\begin{tabular}{L{1.95cm}L{4.55cm}L{3.15cm}L{1.85cm}L{3.40cm}}
\toprule
\rowcolor{TableHead} \textbf{Task} & \textbf{ขนาดข้อมูล (Train / Test)} & \textbf{ตัวแปรเป้าหมาย} & \textbf{ตัวชี้วัดหลัก} & \textbf{Validation Split} \\
\midrule
Task 1: Binary Cls. & Train $104{,}000 \times 33$ (90,816 มีป้าย) / Test $25{,}889 \times 32$ & \code{is\_cheater} \code{\{0, 1\}} (31 ฟีเจอร์ตัวเลข) & Binary $F_2$ ($\beta=2$) & 5-Fold Stratified K-Fold (แถวมีป้าย) \\
Task 2: Multi-Class & Train $104{,}000 \times 46$ / Test $25{,}889 \times 45$ (31 ตัวเลข + 3 noise + 10 หมวด) & \code{segment} \code{\{0..3\}} (\code{Casual}, \code{Grinder}, \code{Social}, \code{Whale}) & Unweighted Macro $F_1$ & 5-Fold SKF + Cross-Fit Calibration \\
Task 3: Spending & Train $104{,}000 \times 34$ / Test $25{,}889 \times 33$ (32 ฟีเจอร์รวม \code{segment}) & \code{spending\_30d} $\in [0, 300{,}000]$ THB (48.12\% ศูนย์) & NMAE Score ($L_1$-Median) & Strict Nested 5-Fold Outer (\code{y>0}) \\
Task 4: Image Cls. & Train 31,546 / Test 25,889 ภาพ (RGB $230 \times 120$ px) & \code{label} \code{\{0..4\}} (5 เกมหลัก) & Unweighted Macro $F_1$ & 5-Fold SGKF (64-bit dHash 28,088 ฉาก) \\
Task 5: Anomaly & ไม่มี Train / Test $25{,}889 \times 123$ (29 กลุ่ม $\times 4$ คาบ + 4 Flags) & \code{is\_anomaly} \code{\{0, 1\}} (กัน 4 Flags) & $F_3$ Score ($\beta=3$) & Masked Control (22,926) + Held-Out \\
\bottomrule
\end{tabular}}
\end{table}
\subsection{มาตรการป้องกัน Data Leakage และประตูตรวจสอบสถิติ}

เพื่อป้องกัน Overfitting และรับประกันว่าผลปรับปรุงจะถ่ายทอดสู่ Test Set คณะผู้จัดทำบังคับใช้หลักปฏิบัติป้องกัน Leakage 2 ข้อในทุกงาน:
\begin{enumerate}[itemsep=1pt,topsep=1pt]
    \item \textbf{Pure-Train Invariant:} ขั้นตอนเตรียมข้อมูลทั้งหมด (Imputation, Standardization/Quantile Transform, Density Estimation, Calibration, Weights และ Decision Threshold) จะต้อง Fit บน Training Folds เท่านั้น ส่วน \code{test.csv} กันแยกไว้สำหรับ Inference ขั้นสุดท้ายเพียงอย่างเดียว
    \item \textbf{Paired Bootstrap Significance Gate:} ก่อนยอมรับสถาปัตยกรรมใหม่เข้าสู่ชุดคำตอบหลัก โมเดลต้องแสดงผลปรับปรุงที่สม่ำเสมอข้ามทุก Fold และมีขอบล่างของช่วงความเชื่อมั่น 95\% จากการทำ Paired Row Bootstrap (1,000 รอบ) มากกว่าศูนย์เมื่อเทียบกับโมเดลฐานก่อนหน้า
\end{enumerate}

% ==============================================================================
\clearpage
\section{Task 1: Anti-Cheat Pre-Filter (Binary Classification)}

\subsection{คำอธิบายปัญหา โครงสร้างคุณลักษณะ และบริบททางธุรกิจ}
Task 1 เป็นงาน Binary Classification (\code{is\_cheater = 0} ผู้เล่นสุจริต และ \code{1} ผู้เล่นโกง) จากข้อมูลพฤติกรรมการเล่นและประวัติบัญชีรวม 31 Numerical Features (ไม่นับรหัส \code{player\_id}) ตามเอกสาร \code{task1/data/datacard.txt} โดยจำแนกฟีเจอร์ออกเป็น 3 กลุ่มหลัก ดังนี้
{\raggedright
\begin{itemize}
    \item \textbf{ทักษะการยิงและการเคลื่อนไหว (12 ตัวแปร):} ใช้วัดความแม่นยำและกลไกควบคุม ได้แก่ \code{kill\_death\_ratio}, \code{headshot\_percentage}, \code{win\_rate}, \code{accuracy\_score}, \code{kill\_consistency}, \code{reaction\_time\_ms}, \code{weapon\_switch\_speed}, \code{movement\_pattern\_score}, \code{aiming\_smoothness}, \code{spray\_control\_score}, \code{damage\_per\_round} และ \code{crosshair\_placement}
    \item \textbf{ความเข้าใจเชิงกลยุทธ์และการเล่นเป็นทีม (9 ตัวแปร):} สะท้อนการตัดสินใจเชิงยุทธวิธี ได้แก่ \code{game\_sense\_score}, \code{map\_knowledge}, \code{communication\_rate}, \code{team\_play\_score}, \code{buy\_decision\_score}, \code{clutch\_success\_rate}, \code{first\_blood\_rate}, \code{survival\_time\_avg} และ \code{utility\_usage\_rate}
    \item \textbf{ประวัติบัญชีและพฤติกรรมผิดปกติ (10 ตัวแปร):} บ่งชี้ลักษณะบัญชีและความต่อเนื่องในการเข้าเล่น ได้แก่ \code{account\_age\_days}, \code{level}, \code{level\_progression\_speed}, \code{friend\_network\_size}, \code{reports\_received}, \code{device\_changes\_count}, \code{input\_consistency\_score}, \code{avg\_session\_length\_min}, \code{sessions\_per\_day} และ \code{night\_play\_ratio}
\end{itemize}\par}

เป้าหมายเชิงปฏิบัติการภายใต้มาตรวัด Binary $F_2$ Score คือช่วยทีมผู้ตรวจสอบจัดลำดับความเร่งด่วนของผู้เล่นต้องสงสัยจากรายงานกว่า 10,000 กรณีต่อวันให้ตรวจจับผู้ใช้ Subtle Aimbot และ Wallhack ได้มากที่สุดโดยไม่ระงับบัญชีผู้เล่นสุจริตโดยผิดพลาด ทั้งนี้ตามคำแนะนำในโจทย์เรื่องการจัดการค่าว่างและความไม่สมดุล (Class Weights, Resampling หรือ SMOTE) คณะผู้จัดทำทดสอบแล้วพบว่าเทคนิค SMOTE ทำให้ค่าความน่าจะเป็นทำนายคลาดเคลื่อนเมื่อนำไปใช้บน Test Set ที่มีสัดส่วน Missingness Shift จึงเลือกคงความถูกต้องเชิงความน่าจะเป็น (Probability Calibration) ควบคู่กับ Decision Threshold $\tau^*$ ที่ปรับจูนแบบ Out-of-Fold แทน

\subsection{การสำรวจข้อมูลเบื้องต้น (EDA)}
ชุดข้อมูลฝึกสอน \code{train.csv} ประกอบด้วย 104,000 แถว (33 คอลัมน์) และ \code{test.csv} ประกอบด้วย 25,889 แถว (32 คอลัมน์) ในชุดฝึกสอนมีข้อมูลที่มีป้ายกำกับจำนวน \textbf{90,816 แถว (87.32\%)} แบ่งเป็นผู้เล่นสุจริต 59,091 แถว (65.07\%) และผู้เล่นโกง 31,725 แถว (34.93\%) ส่วนอีก \textbf{13,184 แถว (12.68\%)} ไม่มีป้ายกำกับ (\code{NaN}) จากการวิเคราะห์การแจกแจงและอัตราข้อมูลขาดหาย (\textbf{รูปที่~\ref{fig:t1_dist}} และ \textbf{รูปที่~\ref{fig:t1_miss}}) สรุปข้อค้นพบ 3 ประการ

\begin{enumerate}
    \item \textbf{ภาวะถูกตัดขอบขวา (Right-Censoring) และขอบเขตเฉพาะกลุ่มบัญชี (Bounded Cohorts):} ตัวแปร \code{game\_sense\_score}, \code{aiming\_smoothness}, \code{spray\_control\_score} และ \code{win\_rate} มีจุดสะสมหนาแน่นที่ขอบบน (\code{1.0} หรือ \code{100.0}) เนื่องจากภาวะถูกจำกัดขอบบน (Right-Censoring Cap) ขณะที่กลุ่มตัวแปรประวัติบัญชี (\code{account\_age\_days}, \code{level}, \code{level\_progression\_speed}, \code{friend\_network\_size}, \code{reports\_received}, \code{crosshair\_placement}) แสดงการแบ่งช่วงขอบเขตชัดเจนที่แยกบัญชีโกงบางกลุ่มออกจากผู้เล่นปกติ
    \item \textbf{ความต่างของสัดส่วนข้อมูลขาดหายระหว่าง Train และ Test (Missingness Shift):} เช่น \code{account\_age\_days} ขาดหาย 14.44\% ใน Train แต่ขาดหายเพียง 0.56\% ใน Test, \code{avg\_session\_length\_min} ขาดหาย 22.98\% เทียบกับ 1.58\%, ขณะที่ \code{headshot\_percentage} ขาดหายเพียง 2.70\% ใน Train แต่สูงถึง 14.11\% ใน Test และ \code{win\_rate} ขาดหาย 4.98\% เทียบกับ 13.74\% ลักษณะนี้ทำให้โมเดลที่แทนค่าว่างแบบคงที่เสียประสิทธิภาพเมื่อนำไปใช้บน Test Set
    \item \textbf{ความขัดแย้งระหว่างทักษะการยิงกับความเข้าใจเชิงกลยุทธ์:} ผู้เล่นที่ใช้โปรแกรมช่วยเล็งมีค่าความแม่นยำ (\code{kill\_death\_ratio}, \code{headshot\_percentage}) สูงผิดปกติเมื่อเทียบกับระดับความรู้แผนที่และทักษะการตัดสินใจ (\code{game\_sense\_score}, \code{map\_knowledge})
\end{enumerate}

\begin{figure}[H]\centering
\makebox[\textwidth][c]{\includegraphics[width=1.04\textwidth,height=0.72\textheight,keepaspectratio]{assets/task1_feature_distributions.png}}
\caption{การแจกแจงของฟีเจอร์ทั้ง 31 ตัวแปรใน Task 1 แยกตามกลุ่มสุจริตและผู้โกง แสดงการกระจุกตัวที่ขอบเขตบนสูงสุดและช่วงขอบเขตข้อมูลบัญชี}
\label{fig:t1_dist}
\end{figure}

\begin{figure}[H]\centering
\wideimg[1.00\textwidth]{assets/task1_missing_rate_train_test.png}
\caption{การเปรียบเทียบอัตราข้อมูลขาดหาย (Missing Rate) ระหว่าง Train Set และ Test Set ใน Task 1}
\label{fig:t1_miss}
\end{figure}

\subsection{การเตรียมข้อมูลและการสร้างคุณลักษณะ}
ดำเนินการภายใต้ 5-Fold Stratified K-Fold (\code{SEED=42}) บนข้อมูลที่มีป้ายกำกับ 90,816 แถว โดยสร้างคุณลักษณะ 3 ระดับในทุก Training Fold:
\begin{enumerate}
    \item \textbf{Physical Account Invariants:} เติมค่าอายุบัญชีด้วยสมการเชิงฟิสิกส์ของความก้าวหน้า (\code{age\_filled = 0.615 * level / level\_progression\_speed}) พร้อมสร้างอัตราส่วนเชิงพฤติกรรม \code{friends\_per\_day}, \code{reports\_per\_day}, \code{kd\_x\_acc} และ \code{aim\_comp} พร้อมชุดคัดกรอง 28 ฟีเจอร์หลัก (\code{COLS\_20\_FE}: คัดเลือก 20 ฟีเจอร์ตั้งต้นและ 8 ฟีเจอร์สร้างใหม่) ควบคู่กับชุด 31 ฟีเจอร์เดิม
    \item \textbf{Fold-Safe Aim-vs-Tactics Z-Gaps:} คำนวณค่ามาตรฐาน Z-score ของกลุ่มตัวชี้วัดทักษะการยิงและยุทธวิธีจำนวน 11 ตัวแปรภายในแต่ละ Training Fold (คงค่า \code{NaN} เดิมไว้ให้ Tree-Based Models จัดการโดยตรง) แล้วสร้างส่วนต่างเชิงทักษะ ได้แก่ \code{aim\_minus\_tactics\_z}, \code{hs\_minus\_gs\_z} และ \code{xhair\_minus\_map\_z}
    \item \textbf{Analytical 8-Cohort Support \& 5-Profile Tobit-MVN5:} ถอดแบบ Data-Generating Process เชิงวิเคราะห์ผ่าน 2 กลไก คือ (ก) โมเดลความน่าจะเป็นของ 8 Cohorts (4 Clean Cohorts, 4 Cheater Cohorts) โดยอินทิเกรตตัวแปรที่ขาดหายด้วย 12-Point Gauss-Legendre Quadrature และ (ข) แบบจำลอง \textbf{1D Tobit EM} สำหรับ \code{game\_sense\_score} ที่ถูกตัดขอบขวาที่ 1.0 ร่วมกับ Yeo-Johnson Shrunk-Covariance Multivariate Normal EM ($\mathrm{shrink}=0.3$, 12 รอบ EM) เพื่อจำแนกโปรไฟล์ทักษะ 5 กลุ่ม โดยทำ Semi-Supervised Training เฉพาะในแต่ละ Training Fold ร่วมกับข้อมูลไม่มี Label จำนวน 13,184 แถวในไฟล์ Train Set
\end{enumerate}

\subsection{การออกแบบโมเดลและผลลัพธ์การทดลอง}

โมเดลฐาน Logistic Regression ทำคะแนน OOF ROC-AUC $= 0.76782$ และ Nested $F_2 = 0.76637$ เท่านั้น เมื่อเพิ่มความละเอียด Bin Resolution ของ GBDTs (\code{LightGBM max\_bin=2047}, \code{XGBoost hist max\_bin=4096}, \code{CatBoost border\_count=1024}) ร่วมกับกฎสนับสนุนขอบเขตจากชุดฝึกสอน (\code{HardSupport Clamp}: แยกแถวที่มั่นใจสูงได้ \code{det\_y0 = 25,977} ผิด 1 ราย และ \code{det\_y1 = 9,922} ผิด 0 รายบน OOF) แล้วผสม 5 โมเดลด้วย Nested SLSQP Stack ค่า AUC เพิ่มเป็น 0.89477 และ $F_2 = 0.83920$

จากนั้นเมื่อนำมาผสานเข้ากับ Analytical Cohort Logistic Regression และ 5-Profile Semi-Supervised Tobit-MVN5 โดยปรับน้ำหนักตามความครบถ้วนของข้อมูล (\code{(0.70, 0.10, 0.20)} สำหรับแถวครบถ้วน และ \code{(0.45, 0.40, 0.15)} สำหรับแถวที่มีค่าขาดหาย) โมเดลขั้นสุดท้ายบรรลุ \textbf{OOF Binary $F_2 = 0.841552$}, \textbf{ROC-AUC $= 0.896762$}, PR-AUC $= 0.840860$ และ LogLoss $= 0.350709$ ณ จุดตัด $\tau^* = 0.199023$ โดยปรับปรุงคะแนนดีขึ้นครบทั้ง 5 Folds เหนือขั้นตอนก่อนหน้า (\code{0.840965}, Paired Bootstrap 95\% CI \code{[+0.000158, +0.001005]}) รายละเอียดทั้งหมดแสดงใน \textbf{ตารางที่~\ref{tab:t1_ablation}}, \textbf{ตารางที่~\ref{tab:t1_folds}} และ \textbf{รูปที่~\ref{fig:t1_diag}}

\begin{table}[H]\centering
{\tblfont
\caption{เปรียบเทียบผล 5-Fold Out-of-Fold ของแบบจำลองแต่ละขั้นใน Task 1}
\label{tab:t1_ablation}
\begin{tabular}{L{6.8cm}ccccc}
\toprule
\rowcolor{TableHead} \textbf{สถาปัตยกรรม / ขั้นตอนการทดลอง} & \textbf{ROC-AUC} & \textbf{$F_2$ (@0.50)} & \textbf{Nested $F_2$} & \textbf{LogLoss} & \textbf{$\tau^*$} \\
\midrule
1. Logistic Regression Baseline (Fold-Safe) & 0.76782 & 0.48595 & 0.76637 & 0.53577 & 0.1394 \\
2. Full CleanDGP LGBM + HardSupport Clamp & 0.89448 & 0.63491 & 0.83919 & 0.35504 & 0.1681 \\
3. Interior 4-Model Specialist + HardClamp & 0.89400 & 0.63454 & 0.83888 & 0.35574 & 0.1822 \\
4. Semi-Supervised CleanZ LGBM + HardClamp & 0.89417 & 0.63435 & 0.83920 & 0.35539 & 0.1824 \\
5. Semi-Supervised XGBoost 4096 + HardClamp & 0.89391 & 0.62731 & 0.83907 & 0.35583 & 0.1943 \\
6. Semi-Supervised CatBoost 1024 + HardClamp & 0.89166 & 0.63582 & 0.83793 & 0.35833 & 0.1533 \\
7. 5-Stream Pure-Train Nested SLSQP Stack & 0.89477 & 0.63342 & 0.83920 & 0.35466 & 0.1750 \\
8. Analytical Cohort Support + GBDT Hybrid & 0.89526 & 0.63240 & 0.84043 & 0.35257 & 0.1829 \\
9. Cohort LR Account Likelihood Blend & 0.89646 & 0.63199 & 0.84097 & 0.35110 & 0.1935 \\
10. Semi-Sup 5-Profile Tobit-MVN (Standalone) & 0.89664 & 0.63497 & 0.84058 & 0.35052 & 0.1834 \\
\rowcolor{TableHead} \textbf{11. Cohort LR + Tobit-MVN5 Hybrid (Final)} & \textbf{0.89676} & \textbf{0.63281} & \textbf{0.84155} & \textbf{0.35071} & \textbf{0.1990} \\
\bottomrule
\end{tabular}}
\end{table}

\begin{table}[H]\centering
{\tblfont
\caption{คะแนน $F_2$ ราย Fold ของระบบ Hybrid ขั้นสุดท้ายใน Task 1 (ชนะครบทั้ง 5 Folds)}
\label{tab:t1_folds}
\begin{tabular}{cccccc}
\toprule
\rowcolor{TableHead} \textbf{Fold} & \textbf{Cohort + GBDT $F_2$} & \textbf{Cohort LR Blend $F_2$} & \textbf{Final Hybrid $F_2$} & \textbf{$\Delta$ Gain} & \textbf{Threshold $\tau^*$} \\
\midrule
0 & 0.839041 & 0.839418 & 0.840027 & +0.000609 & 0.199023 \\
1 & 0.841513 & 0.841330 & 0.841807 & +0.000477 & 0.199057 \\
2 & 0.841050 & 0.842079 & 0.842689 & +0.000610 & 0.199023 \\
3 & 0.838370 & 0.839360 & 0.839489 & +0.000128 & 0.199023 \\
4 & 0.842151 & 0.842642 & 0.843752 & +0.001110 & 0.199023 \\
\bottomrule
\end{tabular}}
\end{table}

\begin{figure}[H]\centering
\wideimg[1.07\textwidth]{assets/task1_diagnostics.png}
\caption{การวิเคราะห์ผลของ Task 1: เส้นโค้ง ROC ข้ามขั้นตอน การยกระดับ $F_2$ ทั้ง 5 Folds เมทริกซ์ความสับสน OOF และตัวแปรสำคัญ 15 อันดับแรกจาก LightGBM}
\label{fig:t1_diag}
\end{figure}

\begin{figure}[H]\centering
\wideimg[1.06\textwidth]{assets/task1_test_alignment.png}
\caption{การกระจายตัวของความน่าจะเป็นบน Test Set เทียบจุดตัด $\tau^* = 0.1990$ และสัดส่วนการตรวจพบบัญชีโกง (15,270 จาก 25,889 บัญชี = 58.98\%)}
\label{fig:t1_align}
\end{figure}

\subsection{การวิเคราะห์ข้อผิดพลาดและการตีความเชิงธุรกิจ}

จากการวิเคราะห์ผลลัพธ์ระยะเริ่มต้น พบจุดบกพร่องสำคัญ 2 ประการที่มีผลต่อประสิทธิภาพของแบบจำลอง ดังนี้
\begin{itemize}
    \item \textbf{ผลกระทบจาก Default Threshold \code{0.50}:} การใช้ Decision Threshold เริ่มต้นให้คะแนน Default $F_2$ เพียง 0.63281 เทียบกับ 0.84155 เมื่อจูน Threshold ที่เหมาะสม $\tau^* = 0.1990$ ซึ่งสอดคล้องกับเกณฑ์ $F_2$ ที่ให้น้ำหนัก Recall สูงกว่า Precision
    \item \textbf{ความคลาดเคลื่อนจากข้อมูลขาดหายและภาวะตัดขอบขวา:} โมเดลที่ใช้เพียง Tree Ensemble 15 ฟีเจอร์พร้อมการแทนค่าว่างแบบคงที่มีประสิทธิภาพลดลงบนแถวที่มีข้อมูลขาดหายและแถวที่กระจุกตัว ณ ขอบเขตบน (\code{game\_sense\_score = 1.0}) การเสริมสมการ Tobit Likelihood และการถ่วงน้ำหนักแยกตามความครบถ้วนของข้อมูลจึงช่วยยกระดับ $F_2$ ของกลุ่มข้อมูลขาดหายจาก 0.82599 เป็น 0.82776 ขณะที่กลุ่มข้อมูลครบถ้วนได้ 0.86137
\end{itemize}

ในมิติการปฏิบัติงานจริง (\textbf{รูปที่~\ref{fig:t1_align}}) แม้แบบจำลองจะคัดแยกผู้ต้องสงสัย 15,270 / 25,889 บัญชี (58.98\%) เพื่อรักษาค่า Recall ตามวัตถุประสงค์ $F_2$ แต่สามารถคัดกรองบัญชีสุจริตบนขอบเขตแน่นอนออกได้ 8,266 บัญชี (31.93\%) โดยไม่เกิดการระงับบัญชีผิดพลาด และดึงกลุ่มโกงแน่นอน 3,341 บัญชี (12.91\%) เข้าสู่คิวระงับเร่งด่วน ทำให้ทีมผู้ตรวจสอบที่มีขีดจำกัด 500 คดีต่อวันทำงานได้อย่างตรงจุด ลดความสูญเสียรายได้วันละ 15 ล้านบาทและลดปัญหาการระงับบัญชีผู้เล่นสุจริตออกจากระบบ

% ==============================================================================
\clearpage
\section{Task 2: Player Segment Classification (Multi-Class Classification)}

\subsection{คำอธิบายปัญหา โครงสร้างคุณลักษณะ และการสำรวจข้อมูลเบื้องต้น}
Task 2 จัดกลุ่มผู้เล่นออกเป็น 4 กลุ่มพฤติกรรม ได้แก่ \code{0=Casual} (40,957 ราย หรือ 39.38\%), \code{1=Core/Grinder} (26,032 ราย หรือ 25.03\%), \code{2=Social/Hardcore} (21,029 ราย หรือ 20.22\%) และ \code{3=Whale} (15,982 ราย หรือ 15.37\%) จากชุดข้อมูลฝึกสอน $104{,}000 \times 46$ คอลัมน์ และชุดทดสอบ $25{,}889 \times 45$ คอลัมน์ ตามเอกสาร \code{task2/data/datacard.txt} มีฟีเจอร์นำเข้าทั้งหมด 44 ตัวแปร ครอบคลุม 5 กลุ่มข้อมูล ดังนี้
{\raggedright
\begin{itemize}
    \item \textbf{เวลาเล่นและความสม่ำเสมอ (7 ตัวแปร):} \code{play\_frequency}, \code{avg\_session\_duration}, \code{total\_playtime\_hours}, \code{login\_streak}, \code{days\_since\_last\_login}, \code{avg\_match\_length} และ \code{peak\_concurrent\_hours}
    \item \textbf{การใช้จ่ายและการตอบสนองโปรโมชัน (5 ตัวแปร):} \code{total\_spending\_thb}, \code{avg\_monthly\_spending}, \code{spending\_frequency}, \code{vip\_tier} และ \code{responds\_to\_discounts}
    \item \textbf{สังคมและการแข่งขัน (10 ตัวแปร):} \code{friend\_count}, \code{team\_play\_percentage}, \code{chat\_activity\_score}, \code{friend\_invites\_sent}, \code{gifts\_sent\_received}, \code{ranked\_participation\_rate}, \code{tournament\_entries}, \code{competitive\_rank}, \code{win\_rate\_ranked} และ \code{watches\_esports}
    \item \textbf{ความก้าวหน้า ความชอบ และคุณลักษณะบัญชี (19 ตัวแปร):} แบ่งเป็นตัวแปรเชิงปริมาณ 9 ตัวแปร ได้แก่ \code{achievement\_completion\_rate}, \code{collection\_progress}, \code{rare\_items\_count}, \code{speed\_of\_progression}, \code{item\_type\_preference\_cosmetic}, \code{item\_type\_preference\_performance}, \code{item\_type\_preference\_social}, \code{account\_age\_days} และ \code{preferred\_game\_mode} ร่วมกับตัวแปรหมวดหมู่ 10 ตัวแปร ได้แก่ \code{region}, \code{platform}, \code{device\_type}, \code{payment\_method}, \code{language}, \code{account\_status}, \code{player\_type\_tag}, \code{engagement\_level}, \code{loyalty\_tier} และ \code{skill\_tier}
    \item \textbf{ตัวแปรสัญญาณรบกวน (3 ตัวแปร):} \code{random\_metric\_1}, \code{random\_metric\_2} และ \code{random\_metric\_3} ซึ่งไม่มีความสัมพันธ์กับกลุ่มเป้าหมายและถูกคัดออกก่อนสร้างแบบจำลอง
\end{itemize}\par}

จากการสำรวจข้อมูล (EDA) พบว่าชั่วโมงการเล่นสะสมเพียงอย่างเดียวไม่อาจแยกกลุ่ม Grinder, Social และ Whale ออกจากกันได้ จำเป็นต้องพิจารณาสัดส่วนข้ามมิติระหว่างการใช้จ่าย การแข่งขัน และกิจกรรมทางสังคม ภายใต้มาตรวัด \textbf{Unweighted Macro $F_1$} เนื่องจากคลาส 0 มีจำนวนมากกว่าคลาส 3 ถึง 2.56 เท่า การทำนายด้วยค่า \code{argmax} ปกติจึงลดทอนความสำคัญและประสิทธิภาพการจำแนกกลุ่ม Whale และ Hardcore จากนั้นจำเป็นต้องคัดตัวแปร Noise Features \code{random\_metric\_1..3} ออกตั้งแต่ขั้นเตรียมข้อมูลเพื่อไม่ให้ Tree-Based Models แตก Split บนสัญญาณสุ่ม (Overfitting on Noise)

\subsection{การเตรียมข้อมูลและการออกแบบคุณลักษณะ}
หลังจากตัดคอลัมน์รหัสและสัญญาณรบกวน \code{random\_metric\_1..3} ออกทั้งหมด คณะผู้จัดทำแปลงตัวแปรลำดับ (\code{ENG\_MAP: Low=1..Veteran=4}, \code{LOY\_MAP: Bronze=1..Diamond=5}) จนได้ชุดคุณลักษณะสะอาด \code{Clean28} (25 ตัวแปรเชิงปริมาณและ 3 ตัวแปรอันดับ) แล้วสร้างฟีเจอร์เสริมภายในแต่ละ Fold ของ 5-Fold Stratified K-Fold (\code{SEED=42}) ดังนี้
\begin{enumerate}
    \item \textbf{Missing-Aware ZLR Probabilities (\code{Clean28+ZLR}):} มาตรฐานข้อมูล Z-score ร่วมกับตัวแปรบ่งชี้ค่าว่างผ่าน Multinomial Logistic Regression ภายใน Fold เพื่อผลิตความน่าจะเป็นฐาน \code{lr\_p0..lr\_p3}
    \item \textbf{Orthogonal Persona-Axis Z-Scores (56 คอลัมน์):} สังเคราะห์ค่าเฉลี่ยและส่วนเบี่ยงเบนตาม 4 แกนพฤติกรรมหลัก (Spend, Engagement, Social, Competitive) พร้อมความต่างข้ามแกน (เช่น \code{spend\_minus\_engage\_z}, \code{spend\_minus\_comp\_z}) และอัตราส่วนเชิงประสิทธิภาพต่อเวลาเล่น (\code{spend\_per\_hour}, \code{spend\_per\_freq})
\end{enumerate}

\subsection{การพัฒนาแบบจำลองและการปรับเทียบ Affine Bayes Decision Rule}

คณะผู้จัดทำประเมินโมเดลฐาน 7 ตระกูลบน 5-Fold SKF ได้แก่ Parametric Expectation-Maximization Density (\code{ExactParamEM}, OOF LogLoss $= 0.59563$), \code{XGB\_Clean28}, \code{LGB\_Clean28}, \code{CB\_Clean28} (CatBoost), Deep Tabular \code{CleanTabM\_SWA3}, \code{Fact\_Clean28} และ \code{Oven\_Rev} พบว่า \code{ExactParamEM} จับโครงสร้างประชากรได้ดีที่สุด และเมื่อผสานความน่าจะเป็นในสัดส่วน \textbf{0.90 ExactParamEM + 0.10 CB\_Clean28} ช่วยเสริมความแม่นยำบริเวณเส้นแบ่งระหว่างกลุ่ม (Decision Boundary, $\mathrm{margin} < 0.02$)

หลังจากนั้นจึงทำ Probability Calibration ด้วย Zero-Mean Logit Bias 4 พารามิเตอร์ (\code{bias = [+0.002124, +0.002960, -0.006227, +0.001143]}) และประยุกต์ใช้ Decision Rule ที่เหมาะสมที่สุดสำหรับ Macro-$F_1$ ในรูป \textbf{Affine Bayes Fixed-Point Rule}:
\[
\hat{y}_i = \arg\max_{c \in \{0,1,2,3\}} \bigl(\alpha_c p_{i,c} - \beta_c\bigr)
\]
โดย $\alpha = [2.4982, 3.9719, 5.0147, 6.7329]$ และ $\beta = [1.0252, 1.4845, 1.8033, 2.5184]$ ผลลัพธ์จาก Cross-Fitted Validation บรรลุ \textbf{Macro-$F_1 = 0.759799$} (Accuracy 77.21\%, LogLoss 0.595756) ชนะ \code{ExactParamEM} เดี่ยว (0.759201) ครบทั้ง 5 Folds (0.75806, 0.75766, 0.76186, 0.75902, 0.76236) รายละเอียดแสดงใน \textbf{ตารางที่~\ref{tab:t2_ablation}}, \textbf{ตารางที่~\ref{tab:t2_perclass}} และ \textbf{รูปที่~\ref{fig:t2_diag}}

\begin{table}[H]\centering
{\tblfont
\caption{เปรียบเทียบผล 5-Fold Out-of-Fold ของโมเดลทั้ง 7 ตระกูลและขั้นตอนปรับเทียบใน Task 2}
\label{tab:t2_ablation}
\begin{tabular}{L{8.0cm}ccc}
\toprule
\rowcolor{TableHead} \textbf{โมเดล / ขั้นตอนการผสาน} & \textbf{OOF LogLoss} & \textbf{OOF Macro-$F_1$} & \textbf{น้ำหนักผสาน} \\
\midrule
ExactParamEM (Parametric EM Density) & \textbf{0.59563} & 0.75848 (Argmax) & 0.90 \\
XGB\_Clean28 & 0.60190 & 0.75747 & 0.00 \\
LGB\_Clean28 & 0.60208 & 0.75690 & 0.00 \\
CB\_Clean28 & 0.60233 & 0.75705 & 0.10 \\
CleanTabM\_SWA3 & 0.60376 & 0.75589 & 0.00 \\
Fact\_Clean28 & 0.61358 & 0.75571 & 0.00 \\
Oven\_Rev & 0.61651 & 0.75131 & 0.00 \\
Cross-Fitted 7-Model Blend (Raw Argmax) & 0.59576 & 0.75867 & --- \\
ExactParamEM (Cross-Fitted Affine Macro-$F_1$ Rule) & 0.59563 & 0.75920 & 1.00 \\
\rowcolor{TableHead} \textbf{0.90 ExactParamEM + 0.10 CB\_Clean28 + Logit-Bias + Affine Rule} & 0.59576 & \textbf{0.75980} & \textbf{Final} \\
\bottomrule
\end{tabular}}
\end{table}

\begin{table}[H]\centering
{\tblfont
\caption{ประสิทธิภาพรายกลุ่มผู้เล่นของโมเดลขั้นสุดท้าย Task 2 (OOF 104,000 แถว) และสัดส่วนบน Test Set}
\label{tab:t2_perclass}
\begin{tabular}{lccccc}
\toprule
\rowcolor{TableHead} \textbf{กลุ่มผู้เล่น (Segment)} & \textbf{Support} & \textbf{Precision} & \textbf{Recall} & \textbf{$F_1$-Score} & \textbf{Test Predicted} \\
\midrule
0: Casual Player & 40,957 & 0.8082 & 0.8347 & 0.8212 & 10,587 (40.89\%) \\
1: Core / Grinder & 26,032 & 0.7452 & 0.7540 & 0.7496 & 6,582 (25.42\%) \\
2: Social / Hardcore & 21,029 & 0.7286 & 0.7087 & 0.7185 & 5,012 (19.36\%) \\
3: Whale (Big Spender) & 15,982 & 0.7769 & 0.7246 & 0.7498 & 3,708 (14.32\%) \\
\rowcolor{TableHead} \textbf{Macro Average (Unweighted)} & \textbf{104,000} & \textbf{0.7647} & \textbf{0.7555} & \textbf{0.7598} & \textbf{25,889 (100.00\%)} \\
\bottomrule
\end{tabular}}
\end{table}

\begin{figure}[H]\centering
\wideimg[1.07\textwidth]{assets/task2_diagnostics.png}
\caption{การวิเคราะห์ผลของ Task 2: ฟีเจอร์สำคัญของ Clean28+ZLR ค่า OOF LogLoss ของโมเดลทั้ง 7 ตระกูล เมทริกซ์ความสับสน และเสถียรภาพราย Fold}
\label{fig:t2_diag}
\end{figure}

\begin{figure}[H]\centering
\wideimg[1.06\textwidth]{assets/task2_test_alignment.png}
\caption{ความสอดคล้องของสัดส่วนกลุ่มผู้เล่นระหว่าง Train, OOF และ Test Set (25,889 แถว) พร้อม\mbox{การกระจายตัว}ของ Decision Margin}
\label{fig:t2_align}
\end{figure}

\subsection{ความผิดพลาดที่ค้นพบและผลกระทบต่อธุรกิจ}

ข้อผิดพลาดสำคัญที่พบระหว่างการพัฒนาเกิดขึ้นเมื่อโค้ดขั้นตอน Export ละเว้นพจน์ชดเชย $-\beta_c$ ในสมการ $\arg\max_c(\alpha_c p_c - \beta_c)$ ขณะที่การใช้ SLSQP ลด Global LogLoss ตัดน้ำหนัก CatBoost เหลือศูนย์ใน 3 Folds ส่งผลให้ Macro-$F_1$ ลดลงเหลือ 0.753158 (ทำนายคลาดเคลื่อน 1,622 แถว) เมื่อปรับแก้ให้ใช้ฟังก์ชัน \code{predict\_segments} ชุดเดียวกันทั้งใน Validation และ Export พร้อมบังคับใช้น้ำหนัก 0.90/0.10 และ Cross-Fitted Logit-Bias คะแนนเพิ่มขึ้นอย่างมีนัยสำคัญสู่ \textbf{0.759799} (95\% Paired Bootstrap CI เท่ากับ \code{[+0.000159, +0.000805]}) ทั้งนี้ควรระบุเป็นข้อจำกัดทางเทคนิคว่าขั้นตอน Logit Bias Calibration และ Affine Rule ดำเนินการแบบ Cross-Fitted บน OOF Probabilities ของ Frozen Base Models

ในเชิงธุรกิจ (\textbf{รูปที่~\ref{fig:t2_align}}) โมเดลสามารถระบุกลุ่ม Whale ได้ถึง 3,708 บัญชี (14.32\%) และกลุ่ม Social/Hardcore 5,012 บัญชี (19.36\%) ด้วยสัดส่วนที่สอดคล้องกับประชากรจริง ช่วยให้ฝ่ายการตลาดยุติการส่งโปรโมชันส่วนลดไปยังผู้เล่นที่พร้อมจ่ายอยู่แล้ว และจัดสรรงบประมาณ 60 ล้านบาทต่อเดือนได้ตรงกลุ่มเป้าหมาย

% ==============================================================================
\clearpage
\section{Task 3: Monthly Spending Prediction (Regression)}

\subsection{คำอธิบายปัญหา กฎ Bayes-$L_1$ โครงสร้างคุณลักษณะ และการสำรวจข้อมูล}
Task 3 มุ่งพยากรณ์ยอดการใช้จ่ายล่วงหน้า 30 วัน (\code{spending\_30d} หน่วยบาท ในช่วง $[0, 300{,}000]$) ภายใต้มาตรวัด \textbf{NMAE Score} จากชุดข้อมูลฝึกสอน $104{,}000 \times 34$ คอลัมน์ (มีค่าว่างรวม 193,731 ช่อง) และชุดทดสอบ $25{,}889 \times 33$ คอลัมน์ (ค่าว่างรวม 48,266 ช่อง) ตามเอกสาร \code{task3/data/datacard.txt} ตัวแปรนำเข้า 32 คอลัมน์แบ่งออกเป็น 5 กลุ่มหลัก ดังนี้
{\raggedright
\begin{itemize}
    \item \textbf{ประวัติการเงินและโปรโมชัน (11 ตัวแปร):} \code{historical\_spending}, \code{prev\_month\_spending}, \code{total\_transactions}, \code{avg\_transaction\_value}, \code{days\_since\_last\_purchase}, \code{purchase\_frequency}, \code{payment\_methods\_used}, \code{purchases\_on\_discount}, \code{discount\_rate\_used}, \code{seasonal\_spending\_pattern} และ \code{owns\_limited\_edition}
    \item \textbf{สถานะบัญชีและแพลตฟอร์ม (7 ตัวแปร):} \code{account\_age\_days}, \code{vip\_status}, \code{is\_premium\_member}, \code{primary\_game}, \code{games\_played}, \code{cross\_game\_activity} และ \code{platform}
    \item \textbf{กิจกรรมและสังคม (9 ตัวแปร):} \code{friend\_count}, \code{social\_interactions}, \code{guild\_membership}, \code{event\_participation\_rate}, \code{daily\_login\_streak}, \code{avg\_session\_length}, \code{sessions\_per\_week}, \code{total\_playtime\_hours} และ \code{days\_since\_last\_login}
    \item \textbf{ความสำเร็จและการแข่งขัน (4 ตัวแปร):} \code{achievement\_count}, \code{achievement\_completion\_rate}, \code{competitive\_rank} และ \code{tournament\_participation}
    \item \textbf{กลุ่มผู้เล่นที่เชื่อมมาจาก Task 2 (1 ตัวแปร):} \code{segment}
\end{itemize}\par}

จาก EDA (\textbf{รูปที่~\ref{fig:t3_eda}}) พบเหตุผล 2 ประการที่ทำให้การทำ Regression มาตรฐานและการใช้ $\exp(\mathbb{E}[\log(1+Y)])$ โดยตรงไม่ประสบความสำเร็จ (ซึ่งเป็นหลุมพรางหลักของโจทย์ Target Scaling ตามคำแนะนำในเอกสารโครงการ):

\begin{enumerate}
    \item \textbf{ภาวะ Zero-Inflation 48.12\%:} ผู้เล่น 50,049 รายไม่ใช้จ่ายเลย (0.00 บาท) การใช้ค่าคาดหมาย $E[Y \mid X] = p \cdot E[Y \mid Y>0, X]$ จะทำนายค่าบวกขนาดเล็กให้คนกลุ่มนี้เกือบทั้งหมด ทำให้ผลรวม $|y_i - \hat{y}_i|$ สูงเกินความจำเป็น
    \item \textbf{ภาวะ Heavy Right Tail และ Whale Spike:} ในกลุ่มผู้จ่ายเงิน 53,951 ราย (51.88\%) ยอดเงินกระจายตัวกว้าง 5 อันดับขนาด ตั้งแต่ 0.09 บาท, มัธยฐาน 2,314.05 บาท, ค่าเฉลี่ย 20,234.88 บาท, เปอร์เซ็นไทล์ที่ 99 เท่ากับ 184,398.64 บาท จนถึงระดับสูงสุด (Upper Cap) 300,000.00 บาท โดยมีกลุ่มผู้เล่นที่มียอดใช้จ่ายมูลค่าสูง $\mathrm{Uniform}[50{,}000, 300{,}000]$ กระจายตัวอยู่ใน Segment 3 (Whale)
\end{enumerate}

ภายใต้การลดผลรวม $|y_i - \hat{y}_i|$ (สมมูลกับการเพิ่ม NMAE Score) คำตอบ Bayes-$L_1$ Optimal Forecast ที่แท้จริงเมื่อทราบ $p_i = P(Y_i > 0 \mid X_i)$ ตามแนวคิด Two-Stage Hurdle (Stage-1: จะใช้จ่ายหรือไม่, Stage-2: จะจ่ายเท่าใด) คือ:
\[
\hat{y}_i^* = \begin{cases} 0 & \text{ถ้า } p_i \le 0.5 \\ Q_{q^*(p_i)}(Y_i \mid Y_i > 0, X_i), \quad q^*(p_i) = 1 - \dfrac{1}{2 p_i} & \text{ถ้า } p_i > 0.5 \end{cases}
\]

\begin{figure}[H]\centering
\wideimg[1.07\textwidth]{assets/figure1-spending-eda.png}
\caption{การสำรวจข้อมูล Task 3: สัดส่วนผู้ไม่ใช้จ่าย 48.12\% การแจกแจงแบบสองยอดของยอดจ่ายในสเกล $\log_{10}$ ยอดจ่ายราย Segment และความสัมพันธ์กับยอดใช้จ่ายในอดีต}
\label{fig:t3_eda}
\end{figure}

\subsection{การเตรียมข้อมูลและการกู้คืนความสัมพันธ์เชิงพีชคณิต}
ดำเนินการภายใต้ 5-Fold Outer Stratified Split (จัดชั้นบนตัวแปร \code{y > 0}) พร้อม Inner 90/10 Early Stopping:
\begin{enumerate}
    \item \textbf{Deterministic Identity Reconstruction (ความแม่นยำ 99.0\%):} ใช้กฎเอกลักษณ์ทางการเงินที่ค้นพบในข้อมูลเพื่อกู้คืนค่าว่างได้อย่างแม่นยำ ได้แก่ \code{historical\_spending = total\_transactions * avg\_transaction\_value}, \code{purchases\_on\_discount = total\_transactions * discount\_rate\_used} และ \code{is\_premium\_member = (vip\_status >= 3)} ร่วมกับค่ามัธยฐานสัดส่วนราย Fold (\code{prev\_month\_spending / historical\_spending}) ขยายฟีเจอร์สะอาดเป็น 44 คอลัมน์
    \item \textbf{Fold-Safe Segment Recovery:} กู้คืนค่าว่างของ \code{segment} ($\sim 5.5\%$) ภายในแต่ละ Fold ด้วย Smooth-KDE Naive Bayes ร่วมกับ GBDT (\code{LogLoss = 0.1684}, ความแม่นยำราย Fold $\sim 94.2\%$)
    \item \textbf{Multiplicative Spending Law \& Huber Anchor:} บนกลุ่มผู้เล่นที่ยอดใช้จ่ายไม่ติดขอบเขตจำกัดสูงสุด (Uncapped Spenders) ยอดใช้จ่ายเป็นไปตามความสัมพันธ์เชิงคูณใน Log Scale ซึ่งประมาณค่า Baseline ด้วย \code{HuberRegressor} (พารามิเตอร์ \code{epsilon=1.35}, \code{alpha=5.0}, MAE บนกลุ่ม uncapped $= 2{,}723.17$ บาท)
\end{enumerate}

\subsection{วิวัฒนาการจาก Quantile Hurdle สู่ ExactBayesL1V2}

ระบบรุ่นแรก (\code{Quantile Hurdle V3}) ประกอบด้วย Stage-1 จำนวน 6 โมเดลผสม (\code{SLSQP6}) และ Stage-2 พยากรณ์ Quantile Residuals 15 ระดับ (0.02..0.55) จาก 5 ตระกูลโมเดลเพื่อประเมิน $Q_{1 - 1/(2p)}$ ทำคะแนน Strict Outer OOF NMAE ได้ \code{0.795996} (MAE 2,690.26 บาท) สูงกว่าค่ามัธยฐานพื้นฐาน (0.500037, MAE 10,495.48 บาท) และ Log-Linear Ridge Hurdle (0.646699, MAE 5,734.68 บาท)

เพื่อลดความคลาดเคลื่อนจาก Quantile Buckets แบบคงที่ จึงพัฒนาแบบจำลอง \textbf{ExactBayesL1V2} (\textit{Generative Marginalization}) ซึ่งประกอบด้วย (1) ตาราง 32-Cell Exact Activation Table สำหรับ Stage-1 ที่ทำ OOF ROC-AUC $= 0.78535$ (สูงกว่า Logistic Regression 0.77614 และ LightGBM 0.78408), (2) กฎความหนาแน่นเชิงคูณของยอดใช้จ่าย, (3) องค์ประกอบ Segment-3 $\mathrm{Uniform}[50{,}000, 300{,}000]$ และ (4) การแก้สมการหา Conditional Median $F_Y^{-1}(0.5 \mid X)$ โดยตรงด้วย \textbf{28-Step Bisection}

ระบบ \code{ExactBayesL1V2} ทำคะแนน \textbf{Strict 5-Fold Outer OOF NMAE $= 0.796377$} (\textbf{MAE $= 2{,}683.96$ บาท}, เฉลี่ยราย Fold $0.796306 \pm 0.005708$) ชนะ Quantile Hurdle V3 ครบทั้ง 5 Folds (0.786151, 0.798213, 0.799203, 0.798332, 0.799630) ด้วยส่วนต่าง +0.000381 (95\% Paired Bootstrap CI \code{[+0.000212, +0.000544]}, $P = 1.0000$ จาก 1,000 รอบ) และได้รับการคัดเลือกในทุก Nested Inner Fold (\code{w\_v2=1.00, tau=0.500, scale=1.0000}) สรุปใน \textbf{ตารางที่~\ref{tab:t3_ablation}} และ \textbf{รูปที่~\ref{fig:t3_diag}}--\textbf{\ref{fig:t3_test}}

\begin{table}[H]\centering
{\tblfont
\caption{เปรียบเทียบผล 5-Fold Strict Outer OOF Validation ของ Task 3}
\label{tab:t3_ablation}
\begin{tabular}{L{7.4cm}ccc}
\toprule
\rowcolor{TableHead} \textbf{แนวทางการสร้างแบบจำลอง (Paradigm)} & \textbf{OOF MAE (THB)} & \textbf{OOF NMAE} & \textbf{$\Delta$ vs. Median (pp)} \\
\midrule
Fold-Training Median Baseline & 10,495.48 & 0.500037 & 0.0000 \\
5-Fold Log-Linear Ridge Hurdle Baseline & 5,734.68 & 0.646699 & +14.6662 \\
Quantile Hurdle V3 (6-Stage-1 + 5-Family $Q_{\mathrm{FINE}}$) & 2,690.26 & 0.795996 & +29.5959 \\
90\% ExactBayesL1V2 + 10\% Quantile Hurdle V3 Blend & 2,684.15 & 0.796365 & +29.6328 \\
\rowcolor{TableHead} \textbf{ExactBayesL1V2 Generative Marginalization (Final)} & \textbf{2,683.96} & \textbf{0.796377} & \textbf{+29.6340} \\
\bottomrule
\end{tabular}}
\end{table}

\begin{figure}[H]\centering
\wideimg[1.07\textwidth]{assets/figure2-hurdle-calibration-diagnostics.png}
\caption{การวิเคราะห์ผล Stage-1 และ Stage-2 ของ Task 3: Reliability Calibration, Stage-1 ROC (ExactBayesL1V2 32-Cell AUC = 0.78535), Log Parity Plot และ Absolute Error ราย Segment}
\label{fig:t3_diag}
\end{figure}

\begin{figure}[H]\centering
\wideimg[1.06\textwidth]{assets/figure3-residual-error-distribution.png}
\caption{การกระจายตัวของค่าคลาดเคลื่อน OOF ($\pm 15\mathrm{k}$ THB) และฮิสโตแกรม 1,000-Bootstrap Gain เหนือ Quantile Hurdle V3}
\label{fig:t3_res}
\end{figure}

\begin{figure}[H]\centering
\wideimg[1.00\textwidth]{assets/figure4-test-spending-prediction-density.png}
\caption{เปรียบเทียบความหนาแน่นของยอดทำนาย Test Set กับ Train Ground Truth ในสเกล Log (ทำนายศูนย์ 49.6\%, มัธยฐาน 78.76 บาท, สูงสุด 300,000 บาท)}
\label{fig:t3_test}
\end{figure}

\subsection{การวิเคราะห์ข้อผิดพลาดและผลกระทบต่อธุรกิจ}

ดังที่ปรากฏใน \textbf{รูปที่~\ref{fig:t3_diag}(d)} ค่าความคลาดเคลื่อนที่เหลืออยู่ส่วนมากมาจาก \textbf{Segment 3 (Whale)} เนื่องจากลักษณะยอดใช้จ่ายมูลค่าสูงที่กระจายตัวแบบ $\mathrm{Uniform}[50\mathrm{k}, 300\mathrm{k}]$ ไม่สามารถคาดการณ์ค่าจริงล่วงหน้าได้อย่างแม่นยำเกินกว่าค่ามัธยฐานเชิงทฤษฎี ส่วนกลุ่ม Casual, Core และ Hardcore มีความแม่นยำสูงมาก การใช้ Two-Stage Bayes-$L_1$ ยังช่วยแก้ไขข้อบกพร่องของการทำนายด้วย Log-Mean ที่มักให้ค่าสูงเกินจริงและทำให้ผู้เล่นไม่ใช้จ่าย 48.12\% ถูกทำนายเป็นบวก

ในเชิงธุรกิจ ผลพยากรณ์บน Test Set (49.60\% หรือ 12,841 แถวมีค่าเท่ากับ 0.00 บาท, ค่าเฉลี่ย 10,501.85 บาท, มัธยฐาน 78.76 บาท, ยอดรวม 271.88 ล้านบาท) ช่วยให้ฝ่ายการเงินลดความคลาดเคลื่อนการพยากรณ์รายได้รายไตรมาส 35--40\%, ช่วยให้ทีม VIP Support ทั้ง 20 คนจัดลำดับความสำคัญจากมูลค่าใช้จ่ายที่แท้จริงแทนชั่วโมงการเล่น (ป้องกันการสูญเสียลูกค้าระดับสูงอย่างกรณี Whale 255,000 บาท) และยุติการแจกโค้ดส่วนลด 50\% ผิดกลุ่มซึ่งสร้างความเสียหายกว่า 6 ล้านบาทต่อเดือน พร้อมสนับสนุนการวางกำหนดการ Limited-Time Events ได้ตรงช่วงจับจ่ายจริง

% ==============================================================================
\clearpage
\section{Task 4: Game Title Detection (Image Classification)}

\subsection{คำอธิบายปัญหา การป้องกัน Frame Leakage และ EDA}
Task 4 จำแนกภาพหน้าจอการเล่นเกมออกเป็น 5 เกม ได้แก่ \code{0=7-11k} (12,429 ภาพ = 39.40\%), \code{1=fiveN} (4,407 ภาพ = 13.97\%), \code{2=freefried} (5,596 ภาพ = 17.74\%), \code{3=pube} (5,926 ภาพ = 18.79\%) และ \code{4=roblock} (3,188 ภาพ = 10.11\%) รวม 31,546 ภาพใน Train Set และ 25,889 ภาพใน Test Set ตามรายละเอียด \code{task4/data/datacard.txt} ภาพทั้งหมดเป็นไฟล์ JPEG 3 ช่องสี RGB 8-bit/channel ความละเอียด $230 \times 120$ พิกเซล (Crop กลางภาพจากขนาดตั้งต้น $256 \times 144$ หรือ 144p) มาจากแหล่งหลากหลายทั้งสตรีมถ่ายทอดสด กิจกรรมผู้เล่น ฟอรัมชุมชน และ Ticket แจ้งปัญหา จึงมีความแปรผันของแผนที่ ช่วงเวลาในเกม สภาพอากาศ และสถานะหน้าจอ HUD/UI สูง (\textbf{รูปที่~\ref{fig:t4_eda}}) ประเมินด้วย \textbf{Macro $F_1$} โดยโจทย์อนุญาตให้ใช้ Pre-trained Vision Models (เช่น CNN, Vision Transformer, Swin Transformer) บน GPU ได้

จากการทำ EDA พบอุปสรรคสำคัญ 2 ประการที่ส่งผลต่อทั้งการประเมินผลและความทนทานของแบบจำลอง ดังนี้
\begin{itemize}
    \item \textbf{ภาวะ Frame Leakage จากการแบ่งแบบสุ่ม:} ภาพจากฉากเดียวกันหรือเฟรมต่อเนื่องอาจกระจายอยู่ทั้งใน Train และ Validation Fold คณะผู้จัดทำจึงจัดกลุ่มภาพด้วย \textbf{Perceptual 64-bit dHash Clustering} ($\mathrm{Hamming\ distance} \le 10$) ได้ \textbf{28,088 กลุ่มฉากอิสระ} และแบ่ง Fold ด้วย \textbf{5-Fold StratifiedGroupKFold (SGKF)} เพื่อให้การวัดผลสะท้อนภาพที่ไม่เคยเห็นจริง
    \item \textbf{ภาพเสื่อมสภาพ (Corrupted Images) และฉากเดี่ยว:} พบภาพคุณภาพต่ำ/เสียหาย (ความเพี้ยนของสี การเอียง หรือ Noise) จำนวน 10,229 ภาพ ควบคู่กับฉากเดี่ยว 7,931 ภาพ จึงต้องพัฒนาทั้ง Image Decorruption Pipeline (\code{decorrupt\_v4}) และชุดทดสอบความทนทานต่อภาพเสียหายรุนแรง (\code{ExtHard6})
\end{itemize}

\begin{figure}[H]\centering
\wideimg[1.07\textwidth]{assets/figure1-screenshot-eda.png}
\caption{สัดส่วนคลาสทั้ง 5 เกมในชุดฝึกสอน Task 4 (31,546 ภาพ, อัตราส่วนไม่สมดุล 2.04 เท่า) ค่าเฉลี่ยสี RGB และภาพตัวอย่างขนาด $230 \times 120$ พิกเซล}
\label{fig:t4_eda}
\end{figure}

\subsection{สถาปัตยกรรมหลายตระกูลและ Dual-Regime Log-Odds Fusion}

คณะผู้จัดทำพัฒนา Vision Models 2 ตระกูลหลักภายใต้ Strict 5-Fold SGKF ดังนี้:
\begin{enumerate}[leftmargin=1.8em,itemsep=2pt]
    \item \textbf{Scene-Hardened End-to-End Fine-Tuned Backbones:} Train โมเดลแบบ Full-Resolution ($192 \times 384$ พิกเซล) ประกอบด้วย \code{CNXv3-Base}, \code{DINOv3L-Hard} และ \code{EVA02L-Hard} ประเมินทั้งแบบ 2-view ปกติ (\code{Orig2}) และ 6-view ฟื้นฟูภาพ (\code{Hard6 / ExtHard6}) ร่วมกับโมเดล ViT ขนาด 224 พิกเซล (\code{EVA02-Large}, \code{EVA02-Base}, \code{DINOv3-ViT-L}, \code{DINOv3-ViT-B}) และ \code{ConvNeXtV2-StageA/C} / \code{DINOv2-StageA/C}
    \item \textbf{6-Backbone Foundation Engine (\code{FoundVLM-v3-Joint12Blk}):} สกัด Feature Embedding 4 มุมมอง (\code{raw}, \code{raw+hflip}, \code{decorrupt\_v4}, \code{decorrupt\_v4+hflip}) จากโมเดลหลัก 6 รุ่น (\code{SigLIP2-SO400M}, \code{DINOv3-ViT-L}, \code{DINOv2-ViT-G/14}, \code{CLIP-ConvNeXt-Large}, \code{CLIP-ViT-H/14}, \code{DINOv3-ConvNeXt-Base}) ต่อรวมเป็นเวกเตอร์ 13,568 มิติ ปรับมาตรฐานด้วย PowerNorm ($p=0.75$) ผ่าน Dual-Condition \code{CosineResMLP} ที่บังคับความสอดคล้องระหว่างภาพต้นฉบับกับภาพที่ผ่านการฟื้นฟู (Decorrupted) ด้วย Symmetric-KL Consistency Loss
\end{enumerate}

ขั้นตอนการผสานผลลัพธ์ประยุกต์ใช้ \textbf{Nested Log-Odds Geometric Pooling} แบบ Temperature Calibration สองสภาวะควบคู่กับ Shift Gate และ Prior Debiasing ดังนี้
\begin{itemize}
    \item \textbf{Clean Pool (11 Active Streams):} ผสานโมเดลบนภาพหน้าจอสภาวะปกติเพื่อคงความคมชัดขององค์ประกอบ HUD และรายละเอียดระดับสูง
    \item \textbf{Unconstrained \code{ExtHard6} Pool:} กำหนดขอบเขตล่าง \code{min\_bound = 0.00} เพื่อลดน้ำหนัก \code{DINOv2-StageC-6TTA} ที่มี Noise สูงลงเหลือ 0.0000 โดยอัตโนมัติ และเน้นน้ำหนักไปที่ \code{FoundVLM-4Fam-Hard6} (สตรีมของ \code{FoundVLM-v3-Joint12Blk}: 0.2607), \code{CNXv2-StageC-6TTA} (0.2321), \code{EVA02L-Hard-Hard6} (0.2302) และ \code{DINOv3L-Hard-Hard6} (0.2170)
    \item \textbf{Consensus Restoration Shift Gate \& Prior Debiasing:} ควบคุมการสลับระหว่างสองโหมดด้วยเกณฑ์ Jensen-Shannon Divergence $\mathrm{JS}(\mathrm{Orig2}, \mathrm{Hard6}) > \tau_{98\%} = 0.004910$ ร่วมกับเงื่อนไขการลดค่า Predictive Entropy โดยกำหนดน้ำหนักผสานสูงสุด $\gamma = 0.04$ (กระตุ้น Gate บน Test Set จำนวน 795 ภาพ) พร้อมทำ Entropy-Gated Prior Debiasing ($\alpha_{\mathrm{lo}}=0.69 \to \alpha_{\mathrm{hi}}=1.35$)
\end{itemize}

ผลลัพธ์ขั้นสุดท้ายบรรลุ \textbf{Clean Unif-Prior Macro-$F_1 = 0.994636$}, All OOF $F_1 = 0.994054$, Corrupt $F_1 = 0.989205$, \textbf{ExtHard6 $F_1 = 0.969820$} และ OOF LogLoss $= 0.019164$ (\textbf{ตารางที่~\ref{tab:t4_ablation}}, \textbf{ตารางที่~\ref{tab:t4_perclass}} และ \textbf{รูปที่~\ref{fig:t4_diag}}--\textbf{\ref{fig:t4_test}})

\begin{table}[H]\centering
{\tblfont
\caption{เปรียบเทียบผลวิวัฒนาการของโมเดล Task 4 ภายใต้ 5-Fold SGKF (Stage 1--5)}
\label{tab:t4_ablation}
\begin{tabular}{L{7.6cm}ccc}
\toprule
\rowcolor{TableHead} \textbf{โมเดล / ขั้นตอนการพัฒนา} & \textbf{Clean Unif $F_1$} & \textbf{ExtHard6 $F_1$} & \textbf{OOF LogLoss} \\
\midrule
Individual: Base ViT EVA02-Large (224px) & 0.989078 & --- & 0.069633 \\
Individual: Full-Res CNXv3-Base ($192 \times 384$ Hard6) & 0.988908 & 0.950293 & 0.053564 \\
Individual: Full-Res DINOv3L-Hard ($192 \times 384$ Hard6) & 0.991634 & 0.959898 & 0.048022 \\
Individual: Full-Res EVA02L-Hard ($192 \times 384$ Hard6) & 0.991497 & 0.959546 & 0.048044 \\
Individual: FoundVLM-v3-Joint12Blk (6-BB 13,568d) & 0.968200 & 0.925210 & 0.101140 \\
Stage 1: 9-Stream Base (E2E + ViTs + DINOv2) & 0.991313 & --- & 0.026300 \\
Stage 2: Hardened Trio + Restoration Shift Gate & 0.994412 & 0.968800 & 0.019381 \\
Stage 3: + 5-Backbone De-Corrupted FoundVLM-v1 & 0.994567 & 0.968804 & 0.018961 \\
Stage 4: + Dual-Condition SymKL FoundVLM-v2 & 0.994573 & 0.969549 & 0.019192 \\
\rowcolor{TableHead} \textbf{Stage 5 (Final): + Joint-12Blk \& Unconstrained ExtHard6} & \textbf{0.994636} & \textbf{0.969820} & \textbf{0.019164} \\
\bottomrule
\end{tabular}}
\end{table}

\begin{table}[H]\centering
{\tblfont\setlength{\tabcolsep}{3.6pt}
\caption{ประสิทธิภาพรายเกมของโมเดลขั้นสุดท้าย Task 4 และสัดส่วนการทำนายบน Test Set (25,889 ภาพ)}
\label{tab:t4_perclass}
\begin{tabular}{clcccccc}
\toprule
\rowcolor{TableHead} \textbf{ID} & \textbf{ชื่อเกม (Slug)} & \textbf{Support} & \textbf{Clean Prec} & \textbf{Clean Rec} & \textbf{Clean $F_1$} & \textbf{ExtHard6 $F_1$} & \textbf{Test Predicted} \\
\midrule
0 & 7-11 Knights (\code{7-11k}) & 12,429 & 0.9970 & 0.9961 & 0.9966 & 0.9830 & 3,575 (13.81\%) \\
1 & FiveN (\code{fiveN}) & 4,407 & 0.9961 & 0.9930 & 0.9945 & 0.9737 & 6,790 (26.23\%) \\
2 & Free Fried (\code{freefried}) & 5,596 & 0.9957 & 0.9950 & 0.9954 & 0.9722 & 3,932 (15.19\%) \\
3 & PUBE (\code{pube}) & 5,926 & 0.9959 & 0.9941 & 0.9950 & 0.9704 & 6,278 (24.25\%) \\
4 & Roblock (\code{roblock}) & 3,188 & 0.9827 & 0.9950 & 0.9888 & 0.9549 & 5,314 (20.53\%) \\
\bottomrule
\end{tabular}}
\end{table}

\begin{figure}[H]\centering
\makebox[\textwidth][c]{\includegraphics[trim=0 0 486.96bp 0,clip,width=1.02\textwidth]{assets/figure2-validation-diagnostics.png}}\\[0.2cm]
\makebox[\textwidth][c]{\includegraphics[trim=828.24bp 0 0 0,clip,width=0.60\textwidth]{assets/figure2-validation-diagnostics.png}}
\caption{ผลการประเมิน 5-Fold SGKF ของ Task 4: (บน) เมทริกซ์ความสับสน (Accuracy 99.50\%) และค่า $F_1$ รายเกมในโหมดปกติกับ ExtHard6 (ล่าง) น้ำหนัก SLSQP ในกลุ่ม ExtHard6}
\label{fig:t4_diag}
\end{figure}

\begin{figure}[H]\centering
\wideimg[1.06\textwidth]{assets/figure3-regime-error-analysis.png}
\caption{พัฒนาการความแม่นยำ Stage 1--5 ของ Task 4 และการกระจายตัวของ JS Divergence ใน Shift Gate ($\gamma = 0.04$, กระตุ้น Gate 795 ภาพ)}
\label{fig:t4_regime}
\end{figure}

\begin{figure}[H]\centering
\wideimg[0.94\textwidth]{assets/figure4-test-class-distribution.png}
\caption{เปรียบเทียบสัดส่วนคลาสทั้ง 5 เกมระหว่าง Train Set กับผลการทำนายบน Test Set (25,889 ภาพ)}
\label{fig:t4_test}
\end{figure}

\subsection{การวิเคราะห์ข้อผิดพลาดและผลกระทบต่อธุรกิจ}

ระหว่างพัฒนาพบกรณีศึกษาที่ส่งผลลบต่อการจำแนก 2 ประการ ดังนี้
\begin{itemize}
    \item \textbf{การสูญเสียรายละเอียดจากโมเดลความละเอียดต่ำ:} โมเดลความละเอียดต่ำ (\code{CNXv2-StageC-2TTA} และ \code{DINOv2-StageC-6TTA}) มักสูญเสียรายละเอียดองค์ประกอบ HUD ที่คมชัด คณะผู้จัดทำจึงนำ \code{CNXv2-StageC-2TTA} ออกจาก Clean Pool เมื่อมีโมเดลความละเอียดเต็ม \code{CNXv3-Base} ($192 \times 384$) และกำหนด \code{min\_bound = 0.00} เพื่อปรับลดน้ำหนัก \code{DINOv2-StageC-6TTA} ลงเป็น 0.0000 ในโหมดภาพเสียหาย
    \item \textbf{เกณฑ์ Shift Gate ที่มีความไวสูงเกินไป (Over-Aggressive Threshold):} การกำหนดน้ำหนักผสาน $\gamma = 0.08$ ทำให้ภาพหน้าจอปกติบางส่วนถูกเบี่ยงไปยังโหมดภาพเสียหายโดยไม่จำเป็น เมื่อปรับเกณฑ์เป็น $\gamma = 0.04$ จึงช่วยรักษาความแม่นยำทั้งภาพสะอาดและภาพเสียหายได้พร้อมกัน
\end{itemize}

จากการเปรียบเทียบใน \textbf{รูปที่~\ref{fig:t4_test}} พบว่า Test Set มีการเปลี่ยนแปลงสัดส่วนเกมจาก Train Set อย่างชัดเจน (เช่น \textit{7-11 Knights} ลดจาก 39.4\% เหลือ 13.8\% ขณะที่ \textit{FiveN} และ \textit{Roblock} เพิ่มเป็น 26.2\% และ 20.5\%) โมเดลที่ผ่าน Calibration จึงสามารถคัดแยก Ticket แจ้งปัญหา 2,000 ใบต่อวัน (ลดเวลาคัดกรอง 10+ นาทีต่อใบลงสู่ระบบอัตโนมัติ) ตรวจพิสูจน์ภาพกิจกรรม ``Play Multiple Games for Rewards'' กว่า 50,000 ภาพได้ครบ 100\% แทนการสุ่มตรวจเพียง 3\% และจัดหมวดหมู่ภาพชุมชนกว่า 100,000 ภาพต่อสัปดาห์ด้วยความแม่นยำเกิน 99\% ในเกือบทุกเกม

% ==============================================================================
\clearpage
\section{Task 5: Account Security Monitoring (Unsupervised Anomaly Detection)}

\subsection{คำอธิบายปัญหา โครงสร้างข้อมูลอนุกรมเวลา และข้อจำกัดของวิธี Unsupervised มาตรฐาน}
Task 5 เป็นงานตรวจจับบัญชีผิดปกติ (\code{is\_anomaly = 1}) ภายใต้มาตรวัด \textbf{$F_3$ Score} ในสภาวะ \textbf{Unsupervised Pure-Test} โดยชุดข้อมูลมีเฉพาะ \code{test.csv} ขนาด $25{,}889 \times 123$ คอลัมน์ ปราศจากป้ายกำกับ ภายในประกอบด้วยตัวแปรนำเข้า 122 ตัวแปร ตามเอกสาร \code{task5/data/datacard.txt} แบ่งออกเป็น 3 ส่วน ดังนี้
{\raggedright
\begin{itemize}
    \item \textbf{อนุกรมเวลาพฤติกรรม 4 ช่วงเวลาต่อเนื่อง (\code{\_1..\_4}, 29 ตระกูล $\times 4 = 116$ ตัวแปร):}
    \begin{itemize}[leftmargin=1.5em,itemsep=1pt,topsep=1pt]
        \item \textbf{การเข้าสู่ระบบและอุปกรณ์ (7 กลุ่ม):} \code{login\_count}, \code{login\_lat}, \code{login\_lon}, \code{location\_changes}, \code{device\_count}, \code{ip\_hash\_entropy} และ \code{login\_time\_variance}
        \item \textbf{พฤติกรรมและทักษะการเล่น (13 กลุ่ม):} \code{session\_duration\_mean}, \code{session\_duration\_std}, \code{play\_pattern\_consistency}, \code{activity\_level}, \code{skill\_rating}, \code{skill\_trajectory}, \code{kda\_mean}, \code{kda\_std}, \code{accuracy\_mean}, \code{accuracy\_std}, \code{reaction\_time\_mean}, \code{reaction\_time\_std} และ \code{performance\_stability}
        \item \textbf{การเงินและเศรษฐกิจในเกม (5 กลุ่ม):} \code{purchase\_frequency}, \code{purchase\_amount\_mean}, \code{purchase\_amount\_std}, \code{refund\_count} และ \code{currency\_transfers}
        \item \textbf{สังคมและกิลด์ (4 กลุ่ม):} \code{friend\_adds}, \code{friend\_removes}, \code{communication\_frequency} และ \code{guild\_changes}
    \end{itemize}
    \item \textbf{ตัวแปรภาพรวมระดับบัญชี (2 ตัวแปร):} \code{account\_age\_days} และ \code{behavior\_stability}
    \item \textbf{Security Flags สำหรับทวนสอบผล (4 ตัวแปร):} \code{password\_changed}, \code{vpn\_usage}, \code{suspicious\_login\_time} และ \code{mass\_item\_sale} ซึ่งทีมผู้พัฒนากันออกจากการ Train และประเมินโมเดลโดยสิ้นเชิง เพื่อใช้ทวนสอบความน่าเชื่อถือในขั้นตอนสุดท้าย
\end{itemize}\par}

ข้อค้นพบสำคัญเมื่อทดสอบตามคำแนะนำในโจทย์ (Isolation Forest, One-Class SVM และ Dense Autoencoder) คืออัลกอริทึม Full-Space Detectors เหล่านี้ประสบปัญหา Feature Dilution อย่างรุนแรง เนื่องจากบัญชีที่ผิดปกติจริงมักแสดงความผิดปกติรุนแรงบนฟีเจอร์เฉพาะทางเพียง 1--4 ตัวแปร จากกว่า 75+ มิติ ขณะที่ฟีเจอร์อื่นทั้งหมดยังอยู่ในเกณฑ์ปกติ ส่งผลให้ Anomaly Score รวมถูกลดทอนสัญญาณโดยมิติปกติจนแยกไม่ออกจากผู้เล่นเดินทางข้ามพื้นที่จริงหรือผู้เล่นที่เว้นช่วงเวลาโจมตีเกิน 24 ชั่วโมง

\subsection{การสร้างฟีเจอร์ Personal Drift และกฎ 8 Archetypes}

คณะผู้จัดทำแปลงอนุกรมเวลา 4 ช่วงเวลาเป็นฟีเจอร์เปรียบเทียบกับพฤติกรรมเดิมของตนเองเพื่อดึงสัญญาณความผิดปกติออกมาให้ชัดเจน ได้แก่ \textbf{Personal Drift (\code{\_\_drift4})} (ช่วงที่ 4 ลบค่ามัธยฐานช่วง 1--3), \textbf{Population Level \& Spread (\code{\_\_level}, \code{\_\_spread})} สำหรับตรวจจับบัญชีอัตโนมัติ (Bot) ที่มีรูปแบบพฤติกรรมคงที่ผิดธรรมชาติ และ \textbf{Geo Displacement (\code{geo\_\_disp4}, \code{geo\_\_disp\_max})} วัดระยะการเปลี่ยนพิกัดล็อกอินโดยรองรับเส้นแบ่งเขตวัน

เมื่อประยุกต์ \textbf{1D Density-Gap CBLOF ($\ge 100\times$ Core Spacing)} และ \textbf{Soft-Max Robust Tail Score} ($\tau > 18.35$ ครอบคลุม 826 บัญชีปลายหาง) ร่วมกับการค้นหาช่องว่างความหนาแน่น พบกลุ่มความผิดปกติที่แยกตัวจากกลุ่มปกติอย่างเด็ดขาด 7 กลุ่มฟีเจอร์เดี่ยว (เช่น \code{behavior\_stability} มี Empty-Gap กว้างกว่ากลุ่มปกติถึง $38{,}792.1\times$) ผสานกับขอบเขต 12 มิติ (12D Hyperrectangle) สำหรับกลุ่มที่ 8 คือ \code{erratic\_combat} (\code{kda\_std} $\in [0.7, 1.0]$, \code{accuracy\_std} $\in [9.0, 12.0]$, \code{performance\_stability} $\in [0.85, 0.95]$ ที่สังเกต $\ge 6$ ช่วงเวลา, \textbf{รูปที่~\ref{fig:t5_box}}) สามารถแยกบัญชีผิดปกติแบบชัดเจนได้ \textbf{2,052 บัญชี (7.93\%)}

\subsection{การกู้คืนบัญชี Unstable Behavior ด้วย ExactBayesUB}

ในกลุ่มบัญชีที่ตัวแปร \code{behavior\_stability} ขาดหาย (\code{NaN} $\sim 5\%$) คณะผู้จัดทำแทนที่ LightGBM/CatBoost ด้วย Analytical Generative Model แบบ Closed-Form (Zero Fitted Parameters) ชื่อ \textbf{ExactBayesUB} บนตัวกรองเบื้องต้น (\code{purchase\_frequency} $\in [3, 6]$, \code{account\_age\_days} $\ge 795$) โดยอินทิเกรต Joint Density ของพฤติกรรมปกติเทียบกับ Unstable เหนือ \code{purchase\_frequency\_1..4}, \code{login\_count\_1..4} (ใช้ 200-Point Gauss-Legendre Quadrature) และ \code{account\_age\_days}

เมื่อทำ Masked Validation บนบัญชีปกติที่ทราบค่าจริง 22,926 ราย (2,393 รายผ่านตัวกรองเบื้องต้น, 246 รายผิดปกติ) \code{ExactBayesUB} ทำคะแนน \textbf{Masked ROC-AUC $= 0.9940$} (0.93656 ภายในตัวกรอง 2,393 แถว เทียบ 0.91251 ของ LightGBM เดิม) และ \textbf{PR-AUC $= 0.7992$} เมื่อเลือกจุดตัดที่เพิ่มค่า Expected-$F_3$ สูงสุดได้จุดตัด \code{0.11513} ซึ่งกู้คืนบัญชีซ่อนเร้นได้เพิ่มอีก \textbf{+25 บัญชี} รวมเป็น \textbf{2,077 / 25,889 บัญชี (8.02\%)} โดยให้ค่า Expected-$F_3 = 0.99835$ ณ จุดกู้คืน 25 บัญชี (ใกล้จุดสูงสุดบนเส้นโค้งที่ \textbf{0.998363} ภายใน $10^{-5}$) สรุปผลใน \textbf{ตารางที่~\ref{tab:t5_summary}} และ \textbf{รูปที่~\ref{fig:t5_diag}}--\textbf{\ref{fig:t5_align}}

\begin{table}[H]\centering
{\tblfont
\caption{สรุปจำนวนบัญชีแต่ละ Archetype และผลทวนสอบบนชุดตัวแปร Security Flags ที่กันไว้สำหรับการทวนสอบของ Task 5}
\label{tab:t5_summary}
\begin{tabular}{L{5.3cm}ccccc}
\toprule
\rowcolor{TableHead} \textbf{กลุ่มพฤติกรรม (Archetype)} & \textbf{บัญชี} & \textbf{สัดส่วน} & \textbf{เปลี่ยนรหัสผ่าน} & \textbf{ใช้ VPN} & \textbf{ขายไอเทมมาก} \\
\midrule
\code{normal} (บัญชีปกติ) & 23,812 & 91.98\% & 4.90\% & 15.26\% & 2.07\% \\
\code{account\_takeover} & 508 & 1.96\% & \textbf{61.49\%} & \textbf{50.61\%} & \textbf{43.54\%} \\
\code{account\_sale\_boosting} & 310 & 1.20\% & 5.84\% & 14.78\% & 2.40\% \\
\code{low\_activity} & 265 & 1.02\% & 5.67\% & 11.58\% & 1.23\% \\
\code{silent\_account} & 263 & 1.02\% & 5.67\% & 9.09\% & 1.98\% \\
\code{unstable\_behavior} & 246 & 0.95\% & 5.88\% & 14.22\% & 2.98\% \\
\code{erratic\_combat} (12D Box) & 244 & 0.94\% & 3.91\% & 17.09\% & 2.55\% \\
\code{bot\_24\_7} & 207 & 0.80\% & 4.12\% & 18.97\% & 1.01\% \\
\code{unstable\_behavior\_nan} & 25 & 0.10\% & 0.00\% & 8.33\% & \textbf{13.64\%} \\
\code{vpn\_pricing\_abuse} (เฉพาะกลุ่ม) & 9 & 0.03\% & 11.11\% & \textbf{100.00\%} & 0.00\% \\
\bottomrule
\end{tabular}}
\end{table}

\textit{หมายเหตุ:} กฎ \code{vpn\_pricing\_abuse} (\code{purchase\_frequency\_4 > 7.5}) มีบัญชีที่เข้าเงื่อนไขรวม 189 บัญชี แต่ 180 บัญชีซ้อนทับกับ \code{account\_takeover} จึงคงเหลือในตารางลำดับชั้นอิสระไม่ซ้ำซ้อนจำนวน 9 บัญชี

\begin{figure}[H]\centering
\wideimg[1.07\textwidth]{assets/task5_anomaly_detection_diagnostics.png}
\caption{การวิเคราะห์ผลของ Task 5: เกณฑ์แยก Tail Score ที่ $\tau=18.35$ จำนวนบัญชีใน 9 Archetypes การยืนยันผลบน Security Flags และช่องว่างความหนาแน่น (Empty-Gap) ของ \code{behavior\_stability}}
\label{fig:t5_diag}
\end{figure}

\begin{figure}[H]\centering
\wideimg[1.07\textwidth]{assets/erratic-combat-box.png}
\caption{การแยกกลุ่ม \code{erratic\_combat} (244 บัญชี) ภายในขอบเขตทรงสี่เหลี่ยม 12 มิติ (12D Box) ออกจากประชากรปกติอย่างชัดเจน}
\label{fig:t5_box}
\end{figure}

\begin{figure}[H]\centering
\wideimg[1.08\textwidth]{assets/task5_bayes_calibration.png}
\caption{เส้นโค้งเลือกจุดตัด Expected-$F_3$ ของ ExactBayesUB (+25 บัญชี ณ จุดตัด 0.11513) และกราฟ Calibration เทียบประชากรควบคุม 22,926 บัญชี}
\label{fig:t5_cal}
\end{figure}

\begin{figure}[H]\centering
\wideimg[0.84\textwidth]{assets/task5_test_alignment.png}
\caption{สัดส่วนผลลัพธ์การตรวจจับความผิดปกติขั้นสุดท้ายของ Task 5 (ปกติ 23,812 บัญชี = 91.98\%, ผิดปกติ 2,077 บัญชี = 8.02\%)}
\label{fig:t5_align}
\end{figure}

\subsection{การวิเคราะห์ข้อผิดพลาดและผลกระทบต่อธุรกิจ}

ข้อผิดพลาดสำคัญที่พบในระยะเริ่มต้นคือ กฎคัดแยกช่วงหาง (Tail Rule) เดิมจัดกลุ่มบัญชีปกติผิดพลาด 51 รายในช่วง $[10.51, 11.16]$ ของฟีเจอร์ \code{session\_duration\_std\_\_drift4} ซึ่งจากการตรวจสอบทางสถิติพบว่าเป็นเพียงส่วนหางปกติแบบ Symmetric $\pm 3.2\sigma$ Gaussian Distribution ของผู้เล่นที่เวลาเล่นผันแปรตามธรรมชาติ การคัดเงื่อนไขนี้ออกช่วยกำจัด False Alarms ได้ทันที ขณะที่การตรวจทวนเทียบกับ Security Flags ที่กันไว้ยืนยันว่ากลุ่ม \code{account\_takeover} เปลี่ยนรหัสผ่านถึง 61.49\% และเทขายไอเทม 43.54\% (เทียบกลุ่มปกติ 4.90\% และ 2.07\%) ส่วนกลุ่ม \code{vpn\_pricing\_abuse} พบการใช้ VPN ถึง 100.00\% ช่วยลดภาระการตรวจสอบลวงของทีมความปลอดภัยได้อย่างเป็นรูปธรรม

% ==============================================================================
\clearpage
\section{การบูรณาการระบบ การทำงานร่วมกัน บทสรุป และข้อเสนอแนะ}

\subsection{การบูรณาการผลทำนายทั้ง 5 งาน โครงสร้าง LTV และผลตรวจรับไฟล์คำตอบ}

ในภาพรวมเชิงสถาปัตยกรรม ผลทำนายจากทั้ง 5 โมเดลถูกส่งต่อเข้าสู่สคริปต์รวมคำตอบกลาง \code{merge.py} ในโฟลเดอร์ \code{submission/} เพื่อรวมคอลัมน์คำตอบตามลำดับ \code{ANS00001}--\code{ANS25889} และตรวจสอบความถูกต้อง 4 ด้านก่อนใช้งานจริง ได้แก่ ขนาดข้อมูล ($25{,}889 \times 6$ คอลัมน์), ความครบถ้วน (\code{0} ค่าว่าง/ค่าอนันต์), การคงลำดับรหัสตาม \code{sample\_submission.csv} และขอบเขตชนิดข้อมูลของทุกงาน

จากนั้นระบบจึงเชื่อมผลลัพธ์เข้าสู่ \textbf{Player Lifetime Value (LTV) Prediction System} ของบริษัท Karena Thailand ดังบทบาทของแต่ละโมเดลต่อไปนี้
\begin{itemize}
    \item \textbf{Risk-Adjusted LTV (Task 1 และ Task 5):} ใช้ความน่าจะเป็นการโกงและความผิดปกติทางความปลอดภัยเพื่อปรับลดความเสี่ยงและคัดกรองรายได้ปลอมในระบบ
    \item \textbf{การประเมินมูลค่าและจัดสรรทรัพยากร (Task 2 และ Task 3):} ใช้กลุ่มผู้เล่นและยอดใช้จ่ายคาดการณ์ล่วงหน้า 30 วันเพื่อคำนวณมูลค่าคาดหวังและจัดลำดับคิวบริการ VIP
    \item \textbf{การเชื่อมโยงพฤติกรรมข้ามเกม (Task 4):} ระบุชื่อเกมจากภาพหน้าจอเพื่อเชื่อมโยงข้อมูลสู่ฝ่ายบริการลูกค้าและการตลาด
\end{itemize}
สรุปผลตรวจสอบภายใน (Local CV) และผลตรวจสอบไฟล์คำตอบสุดท้ายแสดงใน \textbf{ตารางที่~\ref{tab:overall_results}} และ \textbf{ตารางที่~\ref{tab:submission_audit}}

\begin{table}[H]\centering
{\tblfont
\caption{สรุปสถาปัตยกรรมและผลลัพธ์ Out-of-Fold Cross-Validation ของทั้ง 5 งาน}
\label{tab:overall_results}
\begin{tabular}{L{2.7cm}L{5.4cm}cL{4.6cm}}
\toprule
\rowcolor{TableHead} \textbf{Task} & \textbf{โมเดลขั้นสุดท้ายที่เลือกใช้} & \textbf{ตัวชี้วัด} & \textbf{ผลตรวจสอบภายใน (Local CV)} \\
\midrule
Task 1: Anti-Cheat & 5-Stream Nested SLSQP + Analytical 8-Cohort LR + 5-Profile Tobit-MVN5 Hybrid & Binary $F_2$ & \textbf{0.841552} OOF $F_2$ (ROC-AUC 0.896762, 5-Fold SKF) \\
Task 2: Segment & 0.90 ExactParamEM + 0.10 CB\_Clean28 + 4-Class Logit Bias + Affine Bayes Rule & Macro $F_1$ & \textbf{0.759799} Cross-Fitted Macro-$F_1$ (LogLoss 0.595756) \\
Task 3: Spending & ExactBayesL1V2 Generative Marginalization + 32-Cell Stage-1 Table & NMAE Score & \textbf{0.796377} Strict Outer OOF NMAE (MAE 2,683.96 THB) \\
Task 4: Game Title & Dual-Regime Log-Odds Fusion + FoundVLM-v3-Joint12Blk + Consensus Shift Gate & Macro $F_1$ & \textbf{0.994636} Clean Unif-Prior $F_1$ (0.969820 ExtHard6) \\
Task 5: Security & 8 Exact Empty-Gap / 12D Box Archetype Rules + ExactBayesUB Analytical Integral & $F_3$ Score & \textbf{0.998363} max Expected-$F_3$ (0.99835 ที่ +25, 2,077 flags) \\
\bottomrule
\end{tabular}}
\end{table}

เมื่อรวมผลทำนายทั้ง 5 งานและปิดการแข่งขันอย่างเป็นทางการ กลุ่ม \textbf{ML004} ได้รับผลคะแนน \textbf{Final Private Leaderboard สูงสุดที่ 0.86680} (Run~27, \code{dc410ae}; โดยผู้จัดแข่งขันประกาศใช้ชุดแถวประเมินเดียวกับ Public Leaderboard จึงยืนยันคะแนนคัดเลือกสุดท้าย ณ 0.86680 ขณะที่ชุดส่งล่าสุด Run~30 PR~\#99, \code{1e248d1} ได้ 0.86678 ซึ่งต่างกันเพียง 10 แถวหรือ 0.04\%) ครอง\textbf{อันดับที่ 4 จากทั้งหมด 22 กลุ่ม} (\textbf{รูปที่~\ref{fig:final_leaderboard}}) ตามหลังอันดับ 1 (\code{ML003}, 0.87351) เพียง 0.00671 และสูงกว่าเกณฑ์มาตรฐานเดียวของการแข่งขัน \textbf{Kaggle Baseline (0.72601)} อยู่ \textbf{+0.14079} รวมถึงเหนือกว่าค่าเฉลี่ยสุดท้ายของทั้ง 22 กลุ่ม ($0.84321 \pm 0.02744$, $t = +0.8598$) อย่างชัดเจน โดยมีผลตรวจรับไฟล์ส่งคำตอบขั้นสุดท้ายดัง \textbf{ตารางที่~\ref{tab:submission_audit}}

\begin{figure}[H]\centering
\wideimg[1.06\textwidth]{assets/kaggle_final_leaderboard.png}
\caption{เปรียบเทียบคะแนน Final Private Leaderboard ทั้ง 22 กลุ่ม (ML001--ML022): กลุ่ม ML004 ได้อันดับที่ 4 (0.86680) สูงกว่าค่าเฉลี่ยรุ่น 0.84321 และ Kaggle Baseline (0.72601) อย่างชัดเจน}
\label{fig:final_leaderboard}
\end{figure}

\begin{table}[H]\centering
{\tblfont\renewcommand{\arraystretch}{1.08}\setlength{\tabcolsep}{4pt}
\caption{ผลตรวจสอบความสมบูรณ์และสถิติเชิงพรรณนาของไฟล์ส่งคำตอบ \code{submission.csv} ($25{,}889 \times 6$)}
\label{tab:submission_audit}
\begin{tabular}{L{1.9cm}cL{3.6cm}L{7.8cm}}
\toprule
\rowcolor{TableHead} \textbf{คอลัมน์} & \textbf{ค่าว่าง} & \textbf{ชนิด / ขอบเขตข้อมูล} & \textbf{การกระจายตัวและค่าสถิติที่ตรวจพบจริงบน Test Set} \\
\midrule
\code{id} & 0 & String (\code{ANS00001}--\code{ANS25889}) & ครบ 25,889 รหัสเรียงตรงตาม \code{sample\_submission.csv} \\
\code{task1} & 0 & Integer \code{\{0, 1\}} & Clean \code{0}: 10,619 (41.02\%), Cheater \code{1}: 15,270 (58.98\%) \\
\code{task2} & 0 & Integer \code{\{0..3\}} & Casual: 10,587 (40.89\%), Core: 6,582 (25.42\%), Social: 5,012 (19.36\%), Whale: 3,708 (14.32\%) \\
\code{task3} & 0 & Float $\ge 0$ (THB) & ศูนย์ 12,841 ราย (49.60\%), มัธยฐาน 78.76 บาท, เฉลี่ย 10,501.85 บาท, สูงสุด 300,000.00 บาท \\
\code{task4} & 0 & Integer \code{\{0..4\}} & \code{7-11k}: 3,575 (13.81\%), \code{fiveN}: 6,790 (26.23\%), \code{freefried}: 3,932 (15.19\%), \code{pube}: 6,278 (24.25\%), \code{roblock}: 5,314 (20.53\%) \\
\code{task5} & 0 & Integer \code{\{0, 1\}} & Normal \code{0}: 23,812 (91.98\%), Anomaly \code{1}: 2,077 (8.02\%) \\
\bottomrule
\end{tabular}}
\end{table}

\subsection{การทำงานร่วมกันในทีม โครงสร้างคลังซอร์สโค้ด และบทเรียนสำคัญ}

กลุ่ม \textbf{ML004} แบ่งโครงสร้างการพัฒนา Source Code ออกเป็นโมดูลอิสระแยกตามโฟลเดอร์งาน (\code{task1}--\code{task5}) ภายใต้ระบบควบคุมเวอร์ชัน GitHub บนคลังจัดเก็บไฟล์กลางของกลุ่ม (\href{https://github.com/beambeambeam/ml00x}{\code{beambeambeam/ml00x}}) ซึ่งรวบรวมซอร์สโค้ด Jupyter Notebooks สคริปต์รวมคำตอบ และเอกสารรายงานฉบับสมบูรณ์ไว้ในคลังเดียวกันที่ \href{https://github.com/beambeambeam/ml00x/tree/main/report}{\code{ml00x/tree/main/report}} เพื่อรองรับการพัฒนา Feature Engineering ทดสอบสมมติฐาน และตรวจสอบคะแนน Validation ผ่าน Pull Request โดยไม่ขัดแย้งกัน พร้อมบังคับใช้มาตรฐาน Pure-Train Invariant, กฎ Test Isolation และการทดสอบ Statistical Significance ร่วมกันทั้งกลุ่ม โดยแบ่งผู้รับผิดชอบหลักรายโมดูลตามสมาชิกทั้ง 5 ท่าน ได้แก่ \textit{\mbox{ศุภวิชญ์ มารยาท}} (\code{67070501045}) รับผิดชอบ \textbf{Task 1 (Anti-Cheat Pre-Filter)}, \textit{\mbox{กันต์ธีร์ ดวงมณี}} (\code{67070501003}) รับผิดชอบ \textbf{Task 2 (Player Segment Classification)}, \textit{\mbox{วิศิษฐ์ สุวรรณเนาว์}} (\code{67070501042}) รับผิดชอบ \textbf{Task 3 (Monthly Spending Prediction)} และประสานงานรวมชุดส่งคำตอบผ่านสคริปต์ \code{merge.py} ในโฟลเดอร์ \code{submission/}, \textit{\mbox{พลวริษฐ์ วัฒนเหมรัตน์}} (\code{67070501067}) รับผิดชอบ \textbf{Task 4 (Game Title Detection)} และ \textit{\mbox{นัธทวัฒน์ ปริมสิริคุณาวุฒิ}} (\code{67070501027}) รับผิดชอบ \textbf{Task 5 (Account Security Monitoring)} พร้อมทั้งร่วมพัฒนาและทวนสอบข้ามโมดูลภายในคลังจัดเก็บไฟล์กลางของกลุ่ม ดังสรุปข้อค้นพบใน \textbf{ตารางที่~\ref{tab:insights_mistakes}} และโครงสร้าง Source Code พร้อมผู้รับผิดชอบหลักใน \textbf{ตารางที่~\ref{tab:repo_structure}}

\begin{table}[H]\centering
{\tblfont\renewcommand{\arraystretch}{1.04}\setlength{\tabcolsep}{4pt}
\caption{สังเคราะห์ข้อค้นพบเชิงลึก (Key Insights) และข้อผิดพลาดที่พบบ่อยทั้ง 5 งาน}
\label{tab:insights_mistakes}
\begin{tabular}{L{1.25cm}L{6.8cm}L{7.0cm}}
\toprule
\rowcolor{TableHead} \textbf{Task} & \textbf{ข้อผิดพลาดที่พบบ่อย (Pitfalls)} & \textbf{แนวทางแก้ไขและข้อค้นพบ (Key Insights)} \\
\midrule
Task 1 & ใช้ Threshold 0.50, เติมค่าว่างคงที่โดยไม่สน Shift และใช้ SMOTE จนค่าความน่าจะเป็นคลาดเคลื่อน & จูน Bayes-$F_2$ Threshold ($\tau^* = 0.1990$) ผสาน 1D Tobit EM + Cohort Likelihood บนค่าขอบบน 1.0 \\
Task 2 & คงตัวแปร Noise \code{random\_metric\_1..3}, ใช้ \code{argmax} ตั้งต้น และขาดพจน์ $-\beta_c$ ในขั้นตอน Export & ประเมิน Fixed-Point Affine Bayes Rule ผ่าน \code{predict\_segments} ชุดเดียวทั้ง CV และ Export \\
Task 3 & คาดหมายค่าเฉลี่ย $E[Y|X]$ หรือแปลงค่า $\exp(\hat{z})$ โดยตรง จนทำให้กลุ่ม 0 บาท 48.12\% คลาดเคลื่อน & แก้ Bayes-$L_1$ Median $Q_{1-1/(2p)}$ ด้วย 28-Step Bisection + กฎเอกลักษณ์การเงิน \\
Task 4 & แบ่ง Fold สุ่มจนเกิด Frame Leakage ระหว่างฉากต่อเนื่อง และละเลยภาพเสียหาย & จัดกลุ่มฉากด้วย 64-bit dHash SGKF (28,088 ฉาก) และสลับสู่ \code{ExtHard6} ผ่าน JS Shift Gate \\
Task 5 & ใช้ Isolation Forest/AE บนทุกมิติจนเกิด Feature Dilution และจับหาง Gaussian ผิด & แยกกลุ่มด้วย 1D Empty-Gap ($38{,}792.1\times$) + 12D Box และคำนวณ \code{ExactBayesUB} บนกลุ่ม NaN \\
\bottomrule
\end{tabular}}
\end{table}

\textbf{บทเรียนสำคัญด้าน Machine Learning Engineering จากการแข่งขัน:}
\begin{enumerate}
    \item \textbf{ความเข้าใจ Data-Generating Process ให้ผลสูงกว่าการจูนพารามิเตอร์แบบสุ่ม:} การจำลอง Tobit (Task 1), สมการคูณและกฎ Bayes-$L_1$ สำหรับยอดใช้จ่าย (Task 3) รวมถึง Empty-Gap บน 4 ช่วงเวลา (Task 5) ยกระดับคะแนนข้ามทุก Fold สูงกว่าการทำ Hyperparameter Tuning ของ Tree-Based Models เพียงอย่างเดียว
    \item \textbf{การออกแบบ Validation Split ต้องเคารพความซ้ำซ้อนของข้อมูล:} การใช้ Perceptual dHash จัดกลุ่มฉากก่อนแบ่ง Fold ใน Task 4 และกฎ Pure-Train Invariant ช่วยป้องกันการประเมินค่าประสิทธิภาพ CV ที่สูงเกินความเป็นจริง
    \item \textbf{Single Source of Truth ระหว่าง Validation และ Inference:} ข้อผิดพลาดจากพจน์ $-\beta_c$ ที่ขาดหายไปใน Task 2 เป็นบทเรียนสำคัญว่าโค้ดคำนวณ Decision Rule ในขั้นตอน Validation และขั้นตอน Export ผลลัพธ์จริงต้องใช้ฟังก์ชันเดียวกันเสมอ
\end{enumerate}

\begin{table}[H]\centering
{\tblfont\renewcommand{\arraystretch}{1.00}\setlength{\tabcolsep}{4pt}
\caption{โครงสร้างไฟล์ Source Code รายงาน และผู้รับผิดชอบหลักในคลังกลาง (\code{beambeambeam/ml00x})}
\label{tab:repo_structure}
\begin{tabular}{L{1.40cm}L{2.55cm}L{5.55cm}L{5.40cm}}
\toprule
\rowcolor{TableHead} \textbf{งาน} & \textbf{ผู้รับผิดชอบหลัก} & \textbf{โมดูลและสคริปต์หลัก (\code{ml00x})} & \textbf{Notebook / เอกสารประกอบ} \\
\midrule
Task 1 & ศุภวิชญ์ (1045) & \code{task1/run\_account\_cohorts.py}, \code{run\_gpu\_task1.py} & \code{task1/notebooks/eda.ipynb}, \code{submit/submit.md} \\
Task 2 & กันต์ธีร์ (1003) & \code{task2/run\_bayes\_export.py}, \code{bayes\_decision.py} & \code{task2/task2\_player\_segment.ipynb} \\
Task 3 & วิศิษฐ์ (1042) & \code{task3/run\_upgraded\_task3.py} & \code{task3/task3\_monthly\_spending.ipynb} \\
Task 4 & พลวริษฐ์ (1067) & \code{task4/scripts/fuse\_t4\_sota\_trio.py}, \code{train\_m5\_foundvlm\_v3\_joint12blk.py} & \code{task4/task4\_game\_title\_detection.ipynb} \\
Task 5 & นัธทวัฒน์ (1027) & \code{task5/archetype\_rules.py}, \code{run\_ensemble\_task5.py} & \code{task5/task5\_account\_security.ipynb} \\
Merge \& Report & วิศิษฐ์ (1042) \& ทีม & \code{submission/merge.py}, \code{submit.md} & \code{report/CPE342\_TermProject\_Report.tex} \\
\bottomrule
\end{tabular}}
\end{table}

\subsection{ข้อจำกัด บทสรุป และแนวทางการพัฒนาต่อยอด}

จากการพัฒนาระบบต้นแบบ คณะผู้จัดทำสรุปข้อจำกัดทางเทคนิคในเวอร์ชันปัจจุบัน 3 ประการ ดังนี้
\begin{itemize}[itemsep=1pt,topsep=2pt]
    \item \textbf{การทำงานแบบ Frozen Base Models ใน Task 2:} ขั้นตอน Calibration ด้วย Zero-Mean Logit Bias และ Affine Bayes Rule ยังคงประมวลผลแบบ Cross-Fitted บน OOF Probabilities ของ Frozen Base Models แทนการ Train โมเดลทั้ง 7 ตระกูลใหม่ทั้งหมดภายใน Nested Loop
    \item \textbf{การประเมินแบบไร้ Label ใน Task 5:} เนื่องจากไม่มี Label ใน Test Set ค่า Expected-$F_3$ จึงประเมินผ่านแบบจำลองเชิงวิเคราะห์และกลุ่มควบคุมที่ทราบ \code{behavior\_stability} ร่วมกับ Security Flags อิสระ
    \item \textbf{ข้อกำหนดทรัพยากรคำนวณใน Task 4:} การ Train Deep Vision Backbones ใหม่จากศูนย์จำเป็นต้องใช้ GPU
\end{itemize}

\textbf{บทสรุปและแนวทางพัฒนาในอนาคต:} ระบบสำหรับ \textbf{Karena Player Intelligence System} ของกลุ่ม ML004 บรรลุเป้าหมายทั้ง 5 งานตามข้อกำหนด CPE 342 เพื่อต่อยอดสู่การใช้งานจริงในระดับ Production ควรพัฒนาเพิ่มเติม 3 ด้าน ได้แก่
\begin{itemize}[itemsep=1pt,topsep=2pt]
    \item \textbf{Nested End-to-End Pipeline (Task 2):} ผสานขั้นตอน Calibration และ Affine Rule เข้าสู่รอบ Training หลักแบบเต็มรูปแบบ
    \item \textbf{Scoring Service แบบ Streaming/API:} รวมผลทำนายทั้ง 5 งานมาคำนวณ \textbf{Risk-Adjusted Player LTV} แบบเรียลไทม์
    \item \textbf{ระบบเฝ้าระวัง Data/Concept Drift:} ติดตามการเปลี่ยนแปลงพฤติกรรมผู้เล่นเมื่ออัปเดตแพตช์เกมและรูปแบบทุจริตใหม่
\end{itemize}

% ==============================================================================
\clearpage
\renewcommand{\refname}{บรรณานุกรม}
{\raggedright
\begin{thebibliography}{11}
\addcontentsline{toc}{section}{บรรณานุกรม}

\bibitem{sklearn} F.~Pedregosa \textit{et al.}, ``Scikit-learn: Machine Learning in Python,'' \textit{Journal of Machine Learning Research}, vol.~12, pp.~2825--2830, 2011.

\bibitem{lightgbm} G.~Ke \textit{et al.}, ``LightGBM: A Highly Efficient Gradient Boosting Decision Tree,'' in \textit{Advances in Neural Information Processing Systems (NeurIPS)}, 2017.

\bibitem{xgboost} T.~Chen and C.~Guestrin, ``XGBoost: A Scalable Tree Boosting System,'' in \textit{Proceedings of the ACM SIGKDD International Conference on Knowledge Discovery and Data Mining}, 2016, pp.~785--794.

\bibitem{catboost} L.~Prokhorenkova, G.~Gusev, A.~Vorobev, A.~V.~Dorogush, and A.~Gulin, ``CatBoost: Unbiased Boosting with Categorical Features,'' in \textit{Advances in Neural Information Processing Systems (NeurIPS)}, 2018.

\bibitem{pytorch} A.~Paszke \textit{et al.}, ``PyTorch: An Imperative Style, High-Performance Deep Learning Library,'' in \textit{Advances in Neural Information Processing Systems (NeurIPS)}, 2019.

\bibitem{scipy} P.~Virtanen \textit{et al.}, ``SciPy 1.0: Fundamental Algorithms for Scientific Computing in Python,'' \textit{Nature Methods}, vol.~17, pp.~261--272, 2020.

\bibitem{optuna} T.~Akiba, S.~Sano, T.~Yanase, T.~Ohta, and M.~Koyama, ``Optuna: A Next-generation Hyperparameter Optimization Framework,'' in \textit{Proceedings of the ACM SIGKDD}, 2019, pp.~2623--2631.

\bibitem{tobin1958} J.~Tobin, ``Estimation of Relationships for Limited Dependent Variables,'' \textit{Econometrica}, vol.~26, no.~1, pp.~24--36, 1958.

\bibitem{guo2017} C.~Guo, G.~Pleiss, Y.~Sun, and K.~Q.~Weinberger, ``On Calibration of Modern Neural Networks,'' in \textit{Proceedings of ICML}, 2017, pp.~1321--1330.

\bibitem{dinov2} M.~Oquab \textit{et al.}, ``DINOv2: Learning Robust Visual Features without Supervision,'' \textit{Transactions on Machine Learning Research (TMLR)}, 2024.

\bibitem{eva02} Y.~Fang \textit{et al.}, ``EVA-02: A Visual Representation for Neon Genesis,'' \textit{Image and Vision Computing}, 2024.

\end{thebibliography}\par}

\end{document}
